\chapter{The genus-$1$ seven faces: elliptic collision residues}
\label{ch:genus1-seven-faces}

\providecommand{\PhiFA}{\Phi^{\mathrm{FA}}}
\providecommand{\SpCh}{\mathrm{Sp}^{\mathrm{ch}}}

\begin{abstract}
\textbf{Type signature.} (Open quadrant; bar/twisting presentation on
$(E_\tau, D, \tau)$; the seven faces sit at level~$2$ for
$\pi_\cA^{(1)}$, level~$4$ for $H_z^{\mathrm{KZB}}$ and the elliptic
Gaudin generator, and level~$5$ for the modular shadow connection in
the $d\tau$-direction; hypothesis package: $\cA$ modular Koszul,
non-critical level for affine Kac--Moody, $k\neq -h^\vee$ throughout
the F5--F7 identifications, and finite-type or weight-completed bar
coalgebra.) Modular invariance in the $d\tau$-direction is the level-$5$
modular reconstruction (trace + clutching on
$\overline{\cM}_{1,n}$); it is not an intrinsic closure property.

At genus~$0$, the collision residue $r_\cA(z)$ has one polar kernel:
the bar-cobar twisting morphism on $\mathbb{P}^1$. Rational functions
on $\mathbb{P}^1$ are fixed by their principal part and the
normalisation at infinity; the seven faces of
Theorem~\ref{thm:hdm-seven-way-master} are seven extraction maps to
that Laurent kernel. At genus~$1$, the elliptic curve $E_\tau$ carries
a period matrix, a prime form, and non-trivial $A$- and $B$-cycle
monodromy. These data separate the single genus-$0$ Laurent kernel into
seven elliptic comparison kernels with the same collision principal
part but different global readouts.

For affine Kac--Moody algebras $\cA = \widehat\fg_k$, the seven
genus-$1$ faces are: the elliptic twisting morphism (F1), the
DNP25 elliptic line-operator $r$-matrix (F2), the elliptic
$\lambda$-bracket (F3), the
Knizhnik--Zamolodchikov--Bernard connection (F4), the
Belavin--Drinfeld elliptic $r$-matrix (F5), the elliptic Sklyanin
bracket (F6), and the elliptic Gaudin Hamiltonians (F7). F4 carries
the $d\tau$-direction on moduli, F5 carries root-lattice
quasi-periodicity, and F7 carries the commuting Hamiltonians obtained
from the classical Yang--Baxter equation.

For class~$M$ algebras (Virasoro, $\cW_N$), the genus-$1$ collision
residue produces higher-order elliptic differential operators governed
by the prime-form kernel and its derivatives. These operators sit
outside the affine Hitchin--KZB $r$-matrix lane because their polar
order comes from the higher Virasoro and $\cW$ OPE poles.
The degeneration $\tau \to i\infty$ recovers the genus-$0$ seven
faces uniformly across all seven frameworks.
\end{abstract}


%% ====================================================================
%% 1. THE GENUS-1 BAR PROPAGATOR
%% ====================================================================

\section{The genus-$1$ bar propagator}
\label{sec:g1sf-propagator}

The genus-$0$ bar propagator $\eta_{12} = d\log(z_1 - z_2)$ is a
meromorphic one-form on $\mathbb{P}^1$ with a simple pole along the
diagonal and no other data. On
$E_\tau = \mathbb{C}/(\mathbb{Z} + \mathbb{Z}\tau)$, this object
cannot exist. Liouville's theorem forbids a doubly periodic
meromorphic function with a single simple pole. The propagator is
therefore forced to become quasi-periodic: the price of replacing
$\mathbb{P}^1$ by $E_\tau$ is a $B$-cycle monodromy that every
subsequent construction must absorb.

\begin{definition}[Genus-$1$ bar propagator]
\label{def:g1sf-bar-propagator}
\index{bar propagator!genus-1|textbf}
The \emph{genus-$1$ bar propagator} on $E_\tau$ is the logarithmic
derivative of the prime form:
\begin{equation}\label{eq:g1sf-propagator}
\eta_{12}^{E_\tau}
\;=\;
d\log E(z_1, z_2)
\;=\;
\xi_\tau(z_1 - z_2)\, d(z_1 - z_2),
\end{equation}
where $E(z_1,z_2) = \theta_1(z_1 - z_2 | \tau)/\theta_1'(0|\tau)$
is the prime form on $E_\tau$
(a section of $K^{-1/2} \boxtimes K^{-1/2}$, per the convention of
\S\ref{rem:g1sf-prime-form-conventions}), $\theta_1$ is the Jacobi
theta function, and
\begin{equation}\label{eq:g1sf-xi-kernel}
\xi_\tau(z)
\;:=\;
\partial_z\log\theta_1(z|\tau)
\;=\;
\frac{\theta_1'(z|\tau)}{\theta_1(z|\tau)}
\;=\;
\zeta_\tau(z)-2\eta_1 z
\end{equation}
is the prime-form logarithmic kernel. Here $\zeta_\tau$ is the
Weierstrass zeta function for periods $(1,\tau)$:
\begin{equation}\label{eq:g1sf-weierstrass-zeta}
\zeta_\tau(z)
\;=\;
\frac{1}{z}
+ \sum_{(m,n) \neq (0,0)}
\Bigl[
 \frac{1}{z - m - n\tau}
 + \frac{1}{m + n\tau}
 + \frac{z}{(m + n\tau)^2}
\Bigr].
\end{equation}
\begin{equation}\label{eq:g1sf-P-kernel}
P_\tau(z)
\;:=\;
-\partial_z\xi_\tau(z)
\;=\;
\wp(z,\tau)+2\eta_1 .
\end{equation}
\end{definition}

Near $z = 0$ the prime-form kernel has the Laurent expansion
$\xi_\tau(z) = 1/z + O(z)$, so locally the genus-$1$ propagator
is indistinguishable from its genus-$0$ ancestor $1/z$. The two
diverge globally. The Weierstrass zeta function has quasi-periods
\begin{equation}\label{eq:g1sf-zeta-quasi}
\zeta_\tau(z + 1) = \zeta_\tau(z) + 2\eta_1,
\qquad
\zeta_\tau(z + \tau) = \zeta_\tau(z) + 2\eta_\tau,
\end{equation}
where $\eta_1 = \zeta_\tau(1/2)$ and
$\eta_\tau = \zeta_\tau(\tau/2)$ are related by Legendre's identity
$\eta_1 \tau - \eta_\tau = \pi i$. Subtracting $2\eta_1z$ gives the
prime-form monodromy
\begin{equation}\label{eq:g1sf-xi-quasi}
\xi_\tau(z+1)=\xi_\tau(z),
\qquad
\xi_\tau(z+\tau)=\xi_\tau(z)-2\pi i .
\end{equation}
This $B$-cycle constant is the genus-$1$ datum that no doubly periodic
replacement can eliminate.

\begin{remark}[Prime form conventions]
\label{rem:g1sf-prime-form-conventions}
The prime form $E(z,w)$ on $E_\tau$ is a section of
$K^{-1/2} \boxtimes K^{-1/2}$, consistent with the
bar-propagator conventions used here. Its logarithmic derivative
$d\log E = \eta_{12}^{E_\tau}$ is a meromorphic logarithmic connection
form on the universal cover with a simple pole at $z_1=z_2$ and
residue~$1$. The Arakelov-normalized single-valued propagator is
obtained by adding the standard non-holomorphic period correction; the
holomorphic collision residue in this chapter uses the prime-form
kernel~$\xi_\tau$. The bar propagator has weight~$1$ in both variables,
regardless of the conformal weight of the fields being sewed.
\end{remark}


%% ====================================================================
%% 2. THE GENUS-1 COLLISION RESIDUE
%% ====================================================================

\section{The genus-$1$ collision residue}
\label{sec:g1sf-collision-residue}

The propagator determines the collision residue. At genus~$0$, the
rational kernel $1/(z_1 - z_2)$ extracts Laurent coefficients; the
collision residue $r_\cA(z)$ is the resulting element of
$\End_\cA(2)[\![z^{-1}]\!]$. At genus~$1$, the elliptic kernel
$\xi_\tau(z_1 - z_2)$ extracts the same principal coefficients through
a prime-form normalization, and the collision residue acquires a
$\tau$-dependence that encodes the modular geometry of $E_\tau$.

The construction is canonical: restrict the bar-intrinsic MC element
$\Theta_\cA$ (Theorem~\ref{thm:mc2-bar-intrinsic}) to the collision
stratum of the genus-$1$ moduli space.

\begin{definition}[Genus-$1$ collision residue]
\label{def:g1sf-collision-residue}
\index{collision residue!genus-1|textbf}
For a modular Koszul chiral algebra $\cA$ on an elliptic curve
$E_\tau$, the \emph{genus-$1$ collision residue} is
\begin{equation}\label{eq:g1sf-collision-residue}
r_\cA^{(1)}(z, \tau)
\;:=\;
\operatorname{Res}^{\mathrm{coll}}_{1,2}(\Theta_\cA)
\bigm|_{E_\tau}
\;\in\;
\End_\cA(2)[\![z^{-1}]\!] \otimes \mathbb{C}[\![\tau]\!].
\end{equation}
The superscript $(1)$ records the genus; the dependence on $\tau$
is through the modulus of $E_\tau$.
\end{definition}

The relation between genus-$0$ and genus-$1$ collision residues is
fixed after the prime-form normalization is fixed. The term $1/z$ is
replaced by $\xi_\tau$; the higher poles are replaced by the
derivatives of $P_\tau=-\partial_z\xi_\tau$. This is a statement about
polar kernels. A holomorphic additive constant would not change the
principal part, so it is removed by the prime-form normalization.

\begin{theorem}[Elliptic regularization of the collision residue;
\ClaimStatusConditional]
\label{thm:g1sf-elliptic-regularization}
\index{collision residue!elliptic regularization|textbf}
\textbf{Type signature.} (Open quadrant; bar/twisting presentation on
$(E_\tau, D, \tau)$; level~$2$ for the prime-form normalised collision
residue; hypothesis package: $\cA$ modular Koszul on $E_\tau$ with
finite-type or weight-completed bar coalgebra; prime-form
normalisation of Definition~\ref{def:g1sf-bar-propagator}.)
Let $\cA$ be a modular Koszul chiral algebra with genus-$0$ collision
residue $r_\cA(z) = \sum_{n \geq 0} c_n/z^{n+1}$, where
$c_n = a_{(n)}b$ are the OPE-mode coefficients of the generating
fields. The genus-$1$ collision residue is
\begin{equation}\label{eq:g1sf-elliptic-expansion}
r_\cA^{(1)}(z, \tau)
\;=\;
c_0 \cdot \xi_\tau(z)
\;+\;
\sum_{n \geq 1}
 \frac{(-1)^{n-1}}{n!}\,
 c_n \cdot P_\tau^{(n-1)}(z),
\end{equation}
where $P_\tau^{(m)} = d^mP_\tau/dz^m$. Since
$P_\tau=\wp+2\eta_1$, one has $P_\tau^{(m)}=\wp^{(m)}$ for
$m\geq1$. In particular, the first three terms are
\begin{equation}\label{eq:g1sf-first-three-terms}
r_\cA^{(1)}(z, \tau)
\;=\;
c_0\, \xi_\tau(z)
\;+\;
c_1\, P_\tau(z)
\;-\;
\tfrac{1}{2}\, c_2\, P_\tau'(z)
\;+\; \cdots
\end{equation}
The elliptic kernels satisfy
$P_\tau^{(m)}(z) = (-1)^m (m+1)!/z^{m+2} + O(1)$ near $z = 0$, so the
Laurent expansion of $r_\cA^{(1)}$ near $z = 0$ agrees with
$r_\cA(z)$ at leading order.
\end{theorem}

\begin{proof}
Each rational function $1/z^{n+1}$ ($n \geq 0$) appearing in the
genus-$0$ collision expansion is the meromorphic function on
$\mathbb{P}^1$ with that pole at $z=0$ and vanishing at infinity. On
$E_\tau$, the prime-form normalization fixes the corresponding polar
kernel. For $n=0$ it is $\xi_\tau(z)$, whose monodromy is
\eqref{eq:g1sf-xi-quasi}. For $n\geq1$ it is
$(-1)^{n-1}P_\tau^{(n-1)}(z)/n!$. The prefactor is fixed by the
Laurent expansion
$P_\tau^{(m)}(z)=(-1)^m(m+1)!/z^{m+2}+O(1)$ at $m=n-1$, which gives
$(-1)^{n-1}P_\tau^{(n-1)}(z)/n! = 1/z^{n+1}+O(z^{-n+1})$.
The substitution into the propagator-extraction formula of
Definition~\ref{def:g1sf-bar-propagator} yields the
expansion~\eqref{eq:g1sf-elliptic-expansion}. The identification of
$c_n$ with the OPE-mode coefficients follows from the genus-$0$
computation of Theorem~\ref{thm:hdm-face-1}, applied with the same
OPE data on the elliptic curve.
\end{proof}

\begin{remark}[Quasi-periodicity dichotomy]
\label{rem:g1sf-dichotomy}
The kernel $P_\tau$ and all its derivatives are doubly periodic.
The quasi-periodicity of $r_\cA^{(1)}(z, \tau)$ is therefore inherited
entirely from the $\xi_\tau(z)$ term:
\begin{equation}\label{eq:g1sf-B-cycle-monodromy}
r_\cA^{(1)}(z + \tau, \tau) - r_\cA^{(1)}(z, \tau)
\;=\;
-2\pi i \cdot c_0.
\end{equation}
When $c_0 = 0$ (the ``elliptic-function'' case: $\beta\gamma$, lattice
VOAs), the collision residue is a genuine meromorphic function on
$E_\tau$. When $c_0 \neq 0$ (the prime-form case: affine
Kac--Moody, Virasoro, $\cW_N$), the collision residue is
quasi-periodic with non-trivial $B$-cycle monodromy. This dichotomy
controls the splitting behaviour of the genus-$1$ Swiss-cheese
structure (Theorem~\ref{thm:b-cycle-quantum-group}).
\end{remark}


%% ====================================================================
%% 3. FACE 1: THE GENUS-1 BAR-COBAR TWISTING MORPHISM
%% ====================================================================

\section{Face 1: the genus-$1$ bar-cobar twisting morphism}
\label{sec:g1sf-face-1}

Face~1 is the ancestor. At genus~$0$, the collision residue equals
the universal twisting morphism $\pi_\cA$ of the bar-cobar adjunction
(Theorem~\ref{thm:hdm-face-1}): the propagator-extraction formula
reads off the twisting morphism from the bar differential through the
rational kernel $1/(z-w)$. At genus~$1$, the bar complex is defined
over $E_\tau$ without modification; only the kernel changes. The
rational kernel $1/(z-w)$ is replaced by $\xi_\tau(z-w)$, and the
twisting morphism inherits the elliptic geometry of the propagator.

\begin{theorem}[Face~1 at genus~$1$: elliptic twisting morphism;
\ClaimStatusConditional]
\label{thm:g1sf-face-1}
\index{twisting morphism!genus-1|textbf}
\textbf{Type signature.} (Open quadrant; bar/twisting presentation on
$(E_\tau, D, \tau)$; level~$2$ in the Beilinson tower; hypothesis
package: $\cA$ modular Koszul on $E_\tau$, finite-type or
weight-completed bar coalgebra, prime-form normalisation of
Definition~\ref{def:g1sf-bar-propagator}.)
The genus-$1$ collision residue equals the genus-$1$ twisting
morphism:
\begin{equation}\label{eq:g1sf-face-1}
r_\cA^{(1)}(z, \tau)
\;=\;
\pi_\cA^{(1)}(z, \tau)
\;:=\;
\operatorname{Res}_{w=z}
\bar{a}(w) \cdot \xi_\tau(z - w)\, d(z-w)
\;\in\;
\mathrm{Tw}(\barBch(\cA)|_{E_\tau},\, \cA|_{E_\tau}).
\end{equation}
The identification is the genus-$1$ specialization of the
bar-intrinsic construction
(Theorem~\textup{\ref{thm:mc2-bar-intrinsic}}).
\end{theorem}

\begin{proof}
The bar-intrinsic MC element $\Theta_\cA = D_\cA - d_0$ is defined
on arbitrary smooth curves $X$
(Theorem~\ref{thm:mc2-bar-intrinsic}). Specializing to
$X = E_\tau$ and restricting to the collision stratum
$\partial\overline{\cM}_{1,2}^{\mathrm{coll}}$ replaces the
genus-$0$ propagator $d\log(z_1 - z_2) = d(z_1-z_2)/(z_1-z_2)$ by
the genus-$1$ propagator
$d\log E(z_1, z_2) = \xi_\tau(z_1-z_2)\, d(z_1-z_2)$. The
residue-extraction procedure of
Proposition~\ref{prop:twisting-morphism-propagator} applies
verbatim, with the substitution $1/(z-w) \mapsto \xi_\tau(z-w)$.
The result is~\eqref{eq:g1sf-face-1}.
\end{proof}

The lesson of Face~1 at genus~$1$ is a clean separation of algebra
from geometry. The algebraic structure of the twisting morphism is
unchanged: the bar complex $\barBch(\cA)$, the deconcatenation
coproduct, and the MC equation $D_\cA^2 = 0$ are all
genus-independent. The geometry enters through a single datum: the
propagator kernel. At genus~$0$, the kernel is $1/(z-w)$. At
genus~$1$, it is $\xi_\tau(z-w)$. At genus~$g \geq 2$, it is
$d\log E(z,w)$ for the prime form on $\Sigma_g$
(\S\ref{subsec:higher-genus}). The bar-intrinsic construction absorbs
the curve's holomorphic collision geometry through this substitution;
the single-valued Arakelov correction belongs to the global connection,
not to the polar residue.


%% ====================================================================
%% 4. FACE 4: THE KZB CONNECTION
%% ====================================================================

\section{Face 4: the Knizhnik--Zamolodchikov--Bernard connection}
\label{sec:g1sf-face-4}

Face~1 identifies the collision residue with the twisting morphism.
Face~4 adds the modulus $\tau$ as an independent variable. At
genus~$0$, Face~4 identifies the collision residue with the
Gaiotto--Zeng commuting sphere Hamiltonians
(Theorem~\ref{thm:hdm-face-4}): these are the $dz$-components of the
KZ connection. At genus~$1$, the moduli space
$\overline{\mathcal{M}}_{1,n}$ has an additional direction: $\tau$
itself. Conformal blocks depend on $\tau$, and this dependence
requires a $d\tau$-component that has no genus-$0$ counterpart. The
result is the Knizhnik--Zamolodchikov--Bernard (KZB) connection on
$\mathrm{Conf}_n(E_\tau)$.

The two components are governed by the prime-form kernel and its
derivative. The $dz$-component involves $\xi_\tau(z_{ij})$. The
$d\tau$-component involves $P_\tau(z_{ij})=-\partial_z\xi_\tau(z_{ij})$,
which is doubly periodic. The passage from $\xi_\tau$ to $P_\tau$ is
the passage from the quasi-periodic connection form to its elliptic
curvature kernel.

\begin{theorem}[Face~4 at genus~$1$: KZB connection;
\ClaimStatusConditional]
\label{thm:g1sf-face-4}
\index{KZB connection!as genus-1 Face 4|textbf}
\index{seven faces!genus-1 Face 4|textbf}
\textbf{Type signature.} (Open quadrant; level~$4$ for the
$dz$-component, level~$5$ for the $d\tau$-component; hypothesis
package: $\cA=\widehat{\fg}_k$ with $k\neq -h^\vee$; KZ/Sugawara
prefix convention of Lemma~\ref{lem:g1sf-casimir-bridge}; the
$d\tau$-component is the level-$5$ modular reconstruction of the
$\mathrm{SL}(2,\Z)$-equivariant shadow on $\overline{\cM}_{1,n}$, not
an intrinsic property of the closed Sugawara extension.)
For $\cA = \widehat{\fg}_k$ at level $k \neq -h^\vee$, the
modular shadow connection $\nabla^{\mathrm{hol}}_{1,n}$ on
$\mathrm{Conf}_n(E_\tau)$ restricts, in the level-prefixed
collision-residue convention of Lemma~\ref{lem:g1sf-casimir-bridge},
to the binary KZB kernel:
\begin{equation}\label{eq:g1sf-kzb-full}
\nabla^{\mathrm{KZB}}
\;=\;
d
\;-\;
\underbrace{
\frac{k}{k + h^\vee}
\sum_{i < j}
\Omega_{ij}\, \xi_\tau(z_{ij})\, dz_{ij}
}_{\text{$dz$-component: genus-$1$ collision residue}}
\;-\;
\underbrace{
\frac{k}{k + h^\vee}
\sum_{i < j}
\Omega_{ij}\, P_\tau(z_{ij})\, d\tau
}_{\text{$d\tau$-component: modular Hamiltonian}},
\end{equation}
On the locus $k\neq0$, dividing out the explicit level prefix gives
the unprefixed Bernard KZB convention of
Definition~\textup{\ref{def:kzb-connection}}; at $k=0$ the two
normalizations are distinct, as in Lemma~\ref{lem:g1sf-casimir-bridge}.
Here $z_{ij}=z_i-z_j$ and $\Omega_{ij}$ is the KZ/Sugawara numerator
acting on the $i$-th and $j$-th tensor factors.
\begin{enumerate}[label=\textup{(\roman*)}]
\item The $dz$-component equals the genus-$1$ collision residue:
\begin{equation}\label{eq:g1sf-kzb-z-component}
H_{z,i}^{\mathrm{KZB}}
\;=\;
\frac{k}{k + h^\vee}
\sum_{j \neq i}
\Omega_{ij}\, \xi_\tau(z_{ij})
\;=\;
\sum_{j \neq i}
\bigl[r_{\widehat{\fg}_k}^{(1)}(z_{ij}, \tau)\bigr]_{ij}.
\end{equation}
This is the genus-$1$ specialization of the Gaiotto--Zeng
Hamiltonians of Face~$4$ \textup{(}Theorem~\textup{\ref{thm:hdm-face-4})}.

\item The $d\tau$-component equals the genus-$1$ modular shadow:
\begin{equation}\label{eq:g1sf-kzb-tau-component}
H_{\tau,i}^{\mathrm{KZB}}
\;=\;
\frac{k}{k + h^\vee}
\sum_{j \neq i}
\Omega_{ij}\, P_\tau(z_{ij})
\;=\;
-\sum_{j \neq i}
\partial_z \bigl[r_{\widehat{\fg}_k}^{(1)}(z_{ij}, \tau)\bigr]_{ij}.
\end{equation}
The modular Hamiltonian is the $z$-derivative of the collision
residue: $P_\tau(z) = -\xi_\tau'(z)$.
\end{enumerate}
\end{theorem}

\begin{proof}
\textbf{Part (i).} The affine Kac--Moody collision residue at
genus~$1$ is
$r_{\widehat{\fg}_k}^{(1)}(z, \tau) = k\,\Omega_{\mathrm{KZ}}\, \xi_\tau(z)/(k + h^\vee)$,
by Theorem~\ref{thm:g1sf-elliptic-regularization} applied to the
KZ/Sugawara-normalized presentation of the genus-$0$ collision residue,
consistent with the KZB connection~\cite{Bernard88} and the
Casimir-convention bridge of Lemma~\ref{lem:g1sf-casimir-bridge}.
% Level prefix k/(k+h^v) mandatory in the KZ/Sugawara-prefixed convention;
% at k=0 residue vanishes in this convention. The k=0, non-vanishing
% KZ-convention object Omega_{KZ}/(h^v z) is a distinct r-matrix (not
% this face's k=0 specialization); see Lemma g1sf-casimir-bridge.
Since affine Kac--Moody has
$p_{\max} = 2$ and $k_{\max} = 1$ (the OPE has poles of order
$2$ and $1$; the bar propagator absorbs one order), the
collision expansion has $c_0 = k\,\Omega_{\mathrm{KZ}}/(k+h^\vee)$
and $c_n = 0$ for $n \geq 1$. The
elliptic regularization~\eqref{eq:g1sf-elliptic-expansion} therefore
gives $r_{\widehat{\fg}_k}^{(1)} = c_0\, \xi_\tau(z)$.
Substituting into the Gaiotto--Zeng formula~\eqref{eq:hdm-face-4}
of Face~4 produces~\eqref{eq:g1sf-kzb-z-component}, which is the
$dz$-component of the level-prefixed KZB kernel.

\textbf{Part (ii).} The $d\tau$-component involves the
modular Hamiltonian $D_i$, whose leading term is
$\sum_{j \neq i}P_\tau(z_{ij})\, \Omega_{ij}/(k + h^\vee)$
(Definition~\ref{def:kzb-connection}). In the level-prefixed
collision convention the same term is multiplied by $k$. The identity
$P_\tau(z) = -\xi_\tau'(z)$ then identifies it with the
$z$-derivative of the prefixed collision residue. In the modular shadow
framework, the $d\tau$-component is the genus-$1$ modular shadow
$\mathrm{Sh}_{1,n}(\Theta_\cA)|_{\tau\text{-direction}}$: the
shadow connection on $\overline{\cM}_{1,n}$ has a component in the
$\tau$-direction controlled by the curvature
$\kappa(\cA) \cdot \omega_1$, and this curvature produces the
$P_\tau$-kernel through the bar propagator's $\tau$-variation.
The identification with the KZB modular
Hamiltonian is the content of
Proposition~\ref{prop:kz-from-shadow} applied at genus~$1$.
\end{proof}

\begin{remark}[The KZB connection as the full genus-$1$ shadow connection]
\label{rem:g1sf-kzb-shadow}
The identification of the KZB connection with the modular shadow
connection $\nabla^{\mathrm{hol}}_{1,n}$ is the genus-$1$ instance
of the comparison in
\S\ref{sec:hdm-hbar-identification}: the shadow connection on
$\overline{\cM}_{g,n}$ restricts to known classical connections at
each genus. At genus~$0$, it gives the KZ connection. At genus~$1$,
it gives the KZB connection. At genus~$\geq 2$, the same construction
is the prime-form shadow connection on configuration spaces of
higher-genus curves.
\end{remark}


%% ====================================================================
%% 5. FACE 5: THE ELLIPTIC r-MATRIX
%% ====================================================================

\section{Face 5: the Belavin--Drinfeld elliptic $r$-matrix}
\label{sec:g1sf-face-5}

Face~4 sees the collision residue through the connection it defines.
Face~5 sees it through the classical Yang--Baxter equation it
satisfies. At genus~$0$, the collision residue for affine Kac--Moody
equals the Drinfeld rational $r$-matrix
$\Omega/((k+h^\vee)z)$
% KZ convention from thm:hdm-face-5; at k=0 gives Omega/(h^v z) != 0.
(Theorem~\ref{thm:hdm-face-5}). At genus~$1$, the rational
$r$-matrix has the Belavin--Drinfeld elliptic lift
\cite{Belavin81,BelavinDrinfeld82}. In the non-dynamical affine
comparison lane, the classical Yang--Baxter equation on $E_\tau$,
the prescribed simple pole, and the root-lattice quasi-periodicity fix
the elliptic kernel up to the standard normalization.

\begin{theorem}[Face~5 at genus~$1$: elliptic $r$-matrix;
\ClaimStatusConditional{} \textup{(}identification with collision residue\textup{)};
\ClaimStatusProvedElsewhere{} \textup{(}classical elliptic $r$-matrix:
Belavin 1981, Belavin--Drinfeld 1982\textup{)}]
\label{thm:g1sf-face-5}
\index{R-matrix@$R$-matrix!elliptic!as genus-1 Face 5|textbf}
\index{Belavin r-matrix@Belavin $r$-matrix!genus-1 identification}
\textbf{Type signature.} (Open quadrant; bar/twisting presentation on
$(E_\tau, D, \tau)$; level~$2$ for the collision residue, level~$4$
for the spectral $r$-matrix it identifies; hypothesis package:
$\cA = \widehat{\fg}_k$ simple, $k\neq -h^\vee$,
KZ/Sugawara level prefix per Lemma~\ref{lem:g1sf-casimir-bridge}.)
For $\cA = \widehat{\fg}_k$ with $\fg$ simple and
$k \neq -h^\vee$, the genus-$1$ collision residue equals the
classical elliptic $r$-matrix up to the standard level normalization:
\begin{equation}\label{eq:g1sf-elliptic-rmatrix}
r_{\widehat{\fg}_k}^{(1)}(z, \tau)
\;=\;
\frac{k}{k + h^\vee}\,
r^{\mathrm{ell}}_\fg(z, \tau),
\end{equation}
where $r^{\mathrm{ell}}_\fg(z, \tau)$ is the Belavin elliptic
$r$-matrix. For $\fg = \mathfrak{sl}_2$:
\begin{equation}\label{eq:g1sf-sl2-elliptic}
r^{\mathrm{ell}}_{\mathfrak{sl}_2}(z, \tau)
\;=\;
\xi_\tau(z) \cdot \frac{H \otimes H}{2}
\;+\;
\phi_+(z, \tau)\, E \otimes F
\;+\;
\phi_-(z, \tau)\, F \otimes E,
\end{equation}
where $\{H, E, F\}$ is the standard basis of $\mathfrak{sl}_2$
and the theta-function ratios are
\begin{equation}\label{eq:g1sf-phi-pm}
\phi_\pm(z, \tau)
\;=\;
\frac{\theta_1'(0|\tau)\, \theta_1(z \pm \tfrac{1}{2}|\tau)}%
{\theta_1(z|\tau)\, \theta_1(\pm\tfrac{1}{2}|\tau)}.
\end{equation}
The elliptic $r$-matrix satisfies the classical Yang--Baxter equation
on $E_\tau$:
\begin{equation}\label{eq:g1sf-elliptic-cybe}
[r_{12}^{\mathrm{ell}}(z_1{-}z_2),\, r_{13}^{\mathrm{ell}}(z_1{-}z_3)]
+ [r_{12}^{\mathrm{ell}}(z_1{-}z_2),\, r_{23}^{\mathrm{ell}}(z_2{-}z_3)]
+ [r_{13}^{\mathrm{ell}}(z_1{-}z_3),\, r_{23}^{\mathrm{ell}}(z_2{-}z_3)]
= 0.
\end{equation}
\end{theorem}

\begin{proof}
For the Cartan channel: the Kac--Moody OPE
$J^a(z)\, J^b(w) \sim k\, \delta^{ab}/(z-w)^2 +
f^{ab}_{\;\;c}\, J^c(w)/(z-w)$ has $c_0 = k\,\Omega/(k + h^\vee)$
(after KZ/Sugawara normalization and the dual Coxeter shift). For
$\fg = \mathfrak{sl}_2$ with $a = b$ in the Cartan subalgebra
$\mathfrak{h}$, the genus-$1$ regularization replaces
$1/z \mapsto \xi_\tau(z)$, producing
$\frac{k}{k+h^\vee}\,\xi_\tau(z) \cdot H \otimes H/2$ in the
Cartan channel, which matches the Cartan-channel term of
$\frac{k}{k+h^\vee}\,r^{\mathrm{ell}}_{\mathfrak{sl}_2}$
in~\eqref{eq:g1sf-sl2-elliptic}.

For the root channels: on $E_\tau$, the chiral fields $J^E$, $J^F$
associated to the positive and negative roots satisfy twisted
boundary conditions under the $B$-cycle monodromy. The
collision residue in the root channels acquires the theta-function
ratios $\phi_\pm(z, \tau)$ of~\eqref{eq:g1sf-phi-pm}, the meromorphic
functions on $\mathbb{C}$ with a simple pole at $z=0$ matching the
genus-$0$ residue $1/z$ and with the $B$-cycle quasi-periodicity
dictated by the root lattice
(Belavin~\cite{Belavin81}).

The classical Yang--Baxter equation~\eqref{eq:g1sf-elliptic-cybe}
follows from the collision depth~$2$ projection of the modular MC
equation at genus~$1$
(Theorem~\ref{thm:collision-depth-2-ybe} applied on $E_\tau$):
the infinitesimal braid relation
$[\Omega_{ij}, \Omega_{ik} + \Omega_{jk}] = 0$ is replaced by its
elliptic counterpart, and the resulting identity among Weierstrass
functions is the classical elliptic YBE.

The overall identification~\eqref{eq:g1sf-elliptic-rmatrix} is the
composite of Theorem~\ref{thm:g1sf-elliptic-regularization}
(elliptic regularization of the collision expansion) with the
explicit Belavin computation~\cite{Belavin81} of
$r^{\mathrm{ell}}_{\mathfrak{sl}_2}$ as the normalized solution of the
classical YBE on $E_\tau$ with a simple pole at $z=0$ and
quasi-periodicity dictated by the root system. The
Belavin--Drinfeld classification~\cite{BelavinDrinfeld82} identifies
this normalized kernel with the elliptic $r$-matrix in the affine
comparison lane.
\end{proof}

\begin{remark}[Higher-rank elliptic $r$-matrices]
\label{rem:g1sf-higher-rank}
For $\fg = \mathfrak{sl}_N$ with $N \geq 3$, the elliptic
$r$-matrix of Belavin is a matrix-valued meromorphic function on
$E_\tau$ with poles at $z = 0$ whose residue is the $\mathfrak{sl}_N$
Casimir. The collision residue produces this $r$-matrix by the same
mechanism: elliptic regularization of the Cartan and root channels,
with theta-function ratios $\phi_\alpha(z, \tau)$ indexed by the
roots $\alpha$ of $\mathfrak{sl}_N$. The
identification~\eqref{eq:g1sf-elliptic-rmatrix} holds for all simple
$\fg$.
\end{remark}

\begin{lemma}[Casimir-convention bridge at genus~$1$;
\ClaimStatusConditional]
\label{lem:g1sf-casimir-bridge}
\index{Casimir!convention bridge, genus-1|textbf}
\index{r-matrix@$r$-matrix!level prefix discipline|textbf}
\textbf{Type signature.} (Open quadrant; level~$2$/$4$ convention bridge
at the binary collision residue; hypothesis package: simple Lie algebra
$\fg$ with dual Coxeter number $h^\vee$;
$r$-matrix endomorphism convention as Drinfeld~\cite{Drinfeld85} or
Etingof--Kazhdan~\cite{EK96}; level $k$ generic, with edge cases at
$k\in\{0,-h^\vee\}$ recorded explicitly.)
Two Casimir normalizations circulate in the seven-face literature and
must be carefully distinguished when the level prefix is audited:
\begin{enumerate}[label=\textup{(\roman*)}]
\item $\Omega_{\mathrm{tr}} := \sum_a t^a \otimes t^a \in \fg \otimes \fg$,
with $\{t^a\}$ orthonormal for the normalized trace form
$\langle X,Y\rangle_{\mathrm{tr}} = -\operatorname{tr}_{\mathrm{adj}}(XY)/(2h^\vee)$.
This Casimir is $k$-independent. It appears in the Drinfeld
spectral $r$-matrix~\cite{Drinfeld85} as $r^{\mathrm{Dr}}_k(z) = k\,\Omega_{\mathrm{tr}}/z$
\textup{(}explicit level prefix\textup{)}.
\item $\Omega_{\mathrm{KZ}} := 2h^\vee\Omega_{\mathrm{tr}}$, the
Sugawara/KZ numerator used in this volume. This is not the intrinsic
dual-basis Casimir for the Killing form. If the Killing form is
$2h^\vee\langle-,-\rangle_{\mathrm{tr}}$, then the Killing dual-basis
Casimir is $\Omega_{\mathrm{tr}}/(2h^\vee)$; the KZ numerator instead
contains the Sugawara rescaling which converts the trace-form current
OPE into the KZ transport coefficient. It appears in the
Etingof--Kazhdan KZ/Kazhdan-normalized $r$-matrix as
$r^{\mathrm{KZ}}_k(z) = \Omega_{\mathrm{KZ}}/((k+h^\vee)z)$
\textup{(}implicit level prefix, absorbed in the denominator\textup{)}.
\end{enumerate}
\textbf{Two presentations, two $r$-matrices.} The Drinfeld and KZ
presentations are \emph{distinct} classical $r$-matrices on
$\fg \otimes \fg$: they coincide $k$-generically on their spectral
behaviour only when the Casimir tensors are rescaled in a
$k$-dependent way. Setting $r^{\mathrm{Dr}}_k(z) = r^{\mathrm{KZ}}_k(z)$
forces $k\,\Omega_{\mathrm{tr}} = \Omega_{\mathrm{KZ}}/(k+h^\vee)$;
combined with the volume convention
$\Omega_{\mathrm{KZ}} = 2h^\vee\Omega_{\mathrm{tr}}$, this reduces to
$k(k+h^\vee) = 2h^\vee$, which is a quadratic condition on
$k$ with two specific solutions, not an identity in $k$.
\emph{The two $r$-matrices are therefore different families,}
distinguished by the level prefix they carry. The landscape census
\textup{(}Chapter~\ref{chap:landscape-census}, class-$L$ row\textup{)}
records the trace-form collision residue
$k\,\Omega_{\mathrm{tr}}/z$ and the KZ-normalized residue
$\Omega_{\mathrm{KZ}}/((k+h^\vee)z)$ as representatives related by
Casimir and coupling rescaling. This is the level-prefix discipline
pitfall: the symbol ``$\Omega$'' in literature is a placeholder whose
meaning depends on whether the Drinfeld or the KZ spectral form is
displayed.

\noindent
\textbf{Edge-case dichotomy.} The two normalizations behave differently
at $k \in \{0, -h^\vee\}$:
\begin{enumerate}[label=\textup{(\alph*)}]
\item \emph{Vanishing level $k = 0$.} In the Drinfeld
normalization, $r^{\mathrm{Dr}}_0(z) = 0$: the entire $r$-matrix
vanishes with the level. The associated Yang--Baxter equation becomes
trivial and carries no Lie-theoretic content. In the KZ
normalization, $r^{\mathrm{KZ}}_0(z) = \Omega_{\mathrm{KZ}}/(h^\vee z)$,
which is non-zero: the Lie bracket of $\fg$ persists and generates
the rational KZ connection at residual level. These are
\emph{distinct $r$-matrices}, not the same $r$-matrix viewed in two
ways; the level prefix has full load-bearing content. In this chapter,
equations~\eqref{eq:g1sf-kzb-full},
\textup{\eqref{eq:g1sf-kzb-z-component}},
\textup{\eqref{eq:g1sf-elliptic-rmatrix}},
\textup{\eqref{eq:g1sf-elliptic-gaudin}},
and~\eqref{eq:g1sf-master-equation} use the \emph{KZ/Sugawara
numerator} convention with explicit $k/(k+h^\vee)$ prefix
\textup{(}matching the Belavin--Drinfeld elliptic $r$-matrix and the
Bernard--Felder KZB connection~\cite{Bernard88}\textup{)}; at $k=0$
the stated collision residue $r^{(1)}_{\widehat{\fg}_k}(z,\tau)
= k\,\Omega_{\mathrm{KZ}}\,\xi_\tau(z)/(k+h^\vee)$ vanishes.
The non-vanishing limit $\Omega_{\mathrm{KZ}}/(h^\vee z)$ stated in
Theorem~\ref{thm:hdm-face-5} and repeated in the genus-$0$
degeneration is the KZ-convention object $r^{\mathrm{KZ}}_0$, not the
$k\to 0$ specialization of the KZ/Sugawara-prefixed face identified
in Theorem~\ref{thm:g1sf-face-5}. The two are not related by any
$k$-independent Casimir rescaling; they are distinct spectral
$r$-matrices that agree only on the quadratic locus
$k(k+h^\vee) = 2h^\vee$ and differ at $k=0$, where the first
vanishes while the second does not.
\item \emph{Critical level $k = -h^\vee$.} Both normalizations
are singular, but in complementary ways. The KZ denominator
$1/(k+h^\vee)$ diverges; the Drinfeld prefactor $k$ stays finite at
$-h^\vee$. The combination $k/(k+h^\vee)$ diverges. The resolution is
that at the critical level the Sugawara construction collapses
(Feigin--Frenkel centre enlarges; Theorem~\ref{thm:hdm-face-5} states
$k\neq -h^\vee$ explicitly) and the collision residue
$r^{(1)}_{\widehat{\fg}_{-h^\vee}}$ is \emph{not defined as a spectral
$r$-matrix} in $\fg \otimes \fg[\![z^{-1}]\!]$: the Sugawara
denominator's zero is the category-theoretic manifestation of the
enlarged centre. The Faces~F5$^{(1)}$--F7$^{(1)}$ identifications are
therefore scoped to $k \neq -h^\vee$ throughout this chapter.
\end{enumerate}
\end{lemma}

\begin{proof}
(i)--(ii): The intrinsic Killing dual-basis Casimir is
$\Omega_{\mathrm{tr}}/(2h^\vee)$ on simple~$\fg$. This follows from
the standard identity
$\kappa_{\mathrm{Kil}}(X,Y) = \operatorname{tr}_{\mathrm{adj}}(\mathrm{ad}_X \mathrm{ad}_Y)
= 2h^\vee\,\langle X,Y\rangle_{\mathrm{tr}}$ on the Cartan of
simple~$\fg$~\cite[\S 1.1]{EK96}, combined with the observation that
the Casimir $\sum_a I^a \otimes I_a$ built from dual bases scales as
the inverse of the form under rescaling. Explicitly: if
$\{t^a\}$ is $\langle\cdot,\cdot\rangle_{\mathrm{tr}}$-orthonormal
then $\{t^a/(2h^\vee)^{1/2}\}$ would be Killing-orthonormal up to
conventions of basis duality; working with the genuine non-degenerate
pairing gives the Killing dual-basis tensor
$\Omega_{\mathrm{Kil}} = \Omega_{\mathrm{tr}}/(2h^\vee)$.

The tensor denoted $\Omega_{\mathrm{KZ}}$ in this volume is different:
it is the Sugawara/KZ numerator
$\Omega_{\mathrm{KZ}}=2h^\vee\Omega_{\mathrm{tr}}$, fixed in
Chapter~\ref{chap:landscape-census} and in the Yangian comparison
chapters. This convention keeps the trace-form collision residue
$k\Omega_{\mathrm{tr}}/z$ separate from the KZ transport coefficient
$\Omega_{\mathrm{KZ}}/((k+h^\vee)z)$.

The intertwining constraint $r^{\mathrm{Dr}}_k = r^{\mathrm{KZ}}_k$
reads $k\,\Omega_{\mathrm{tr}}/z = \Omega_{\mathrm{KZ}}/((k+h^\vee)z)$.
Substituting the KZ/Sugawara convention:
$k\,\Omega_{\mathrm{tr}}/z = 2h^\vee\Omega_{\mathrm{tr}}/((k+h^\vee)z)$,
which reduces to $k(k+h^\vee) = 2h^\vee$, a quadratic equation in
$k$ with two specific solutions. The Drinfeld and KZ $r$-matrices
therefore coincide only at two (generically irrational) levels; they
are otherwise distinct. The symbolic identity
``$k\,\Omega_{\mathrm{tr}} = \Omega/(k+h^\vee)$'' that appears in the
literature is an abuse of notation: a level-dependent rescaling of
``$\Omega$'' is silently performed to force the equation. The
non-trivial mathematical content is the dual-Coxeter shift: the
denominator $k+h^\vee$ is the $\sigma_{\widehat{\fg}}$-symmetric level
(Theorem C complementarity), not a purely algebraic normalization.

(a): The Drinfeld formula $r^{\mathrm{Dr}}_k = k\,\Omega_{\mathrm{tr}}/z$
specialises to $0$ at $k=0$ by direct evaluation. The KZ formula
$r^{\mathrm{KZ}}_k = \Omega_{\mathrm{KZ}}/((k+h^\vee)z)$ specialises
to $\Omega_{\mathrm{KZ}}/(h^\vee z) \neq 0$ at $k=0$. These are
distinct classical $r$-matrices on $\fg \otimes \fg$: the first
satisfies a trivial classical Yang--Baxter equation (since it
vanishes), while the second satisfies the non-trivial infinitesimal
braid relation $[\Omega_{ij}, \Omega_{ik} + \Omega_{jk}] = 0$
evaluated on the KZ/Sugawara numerator~\cite{Bernard88}. This is not
a contradiction: they are two different $r$-matrices corresponding to
two different physical/algebraic regimes. The Drinfeld prefactor
$k$ measures the ``quantum'' level of a Kac--Moody symmetry that
vanishes at $k=0$ (abelian limit of the current algebra); the KZ
denominator $k+h^\vee$ measures the inverse Sugawara parameter that
stays finite and non-zero at $k=0$. The two capture different
invariants of the same underlying chiral algebra.

(b): At $k = -h^\vee$, $r^{\mathrm{KZ}}_{-h^\vee}$ is singular (pole
in $k+h^\vee$). The Sugawara construction
$T^{\mathrm{Sug}} = \sum_a {:}J^a J^a{:}/(2(k+h^\vee))$ is
singular at the same point. The Feigin--Frenkel centre
$\mathfrak{z}(\widehat{\fg}) \subset V_{-h^\vee}(\fg)$ is
infinite-dimensional at $k = -h^\vee$ \textup{(}and trivial for
$k \neq -h^\vee$\textup{)}, providing the categorical resolution of
the Sugawara singularity. The collision residue
$r_\cA^{(1)}(z,\tau)$ as defined by
equation~\eqref{eq:g1sf-collision-residue} still makes sense as a
formal element of $\End_\cA(2)[\![z^{-1}]\!] \otimes
\mathbb{C}[\![\tau]\!]$ at $k=-h^\vee$, but its identification with
the Belavin--Drinfeld elliptic $r$-matrix \textup{(}which requires the
$k/(k+h^\vee)$ prefactor to be finite\textup{)} fails. The chapter's
convention $k \neq -h^\vee$ is therefore a \emph{load-bearing hypothesis}
on the Faces~F5$^{(1)}$--F7$^{(1)}$ identifications, not a
typographical convenience.
\end{proof}

\begin{remark}[Casimir normalization check]
\label{rem:g1sf-casimir-normalization-check}
The Casimir-convention bridge of Lemma~\ref{lem:g1sf-casimir-bridge}
matches the convention now used in
Chapter~\ref{chap:landscape-census} of this volume
\textup{(}landscape census, class~$L$ row\textup{)}: the trace-form
representative $k\Omega_{\mathrm{tr}}/z$ and the KZ representative
$\Omega_{\mathrm{KZ}}/((k+h^\vee)z)$ agree only after the stated
Casimir and coupling rescaling. With fixed Casimir tensors they define
different rational functions. Readers applying the
landscape-census formula numerically must fix the level $k$ and the
Casimir convention simultaneously.

Volume~II \textup{(}$A_\infty$ Chiral Algebras and 3D~HT~QFT,
$\mathsf{SC}^{\mathrm{ch,top}}$ bar differential\textup{)} and
Volume~III \textup{(}CY Categories, K3 Hall--Drinfeld double\textup{)} use the
Drinfeld normalization with explicit level prefix throughout; the
translation between volumes is Lemma~\ref{lem:g1sf-casimir-bridge}.
Cross-volume propagation requires that any statement involving a
numerical value of $r(z)$ at fixed $z = z_0$, $k = k_0$ be tagged with
which Casimir convention is in force at that statement; this chapter
tags all Face~F5$^{(1)}$--F7$^{(1)}$ statements with the
KZ/Sugawara-prefixed convention
$r^{(1)} = k\,\Omega_{\mathrm{KZ}}\xi_\tau/(k+h^\vee)$
\textup{(}explicit $k$ prefix, matching Bernard--Felder KZB\textup{)}.
\end{remark}


%% ====================================================================
%% 6. FACE 7: THE ELLIPTIC GAUDIN MODEL
%% ====================================================================

\section{Face 7: the elliptic Gaudin model}
\label{sec:g1sf-face-7}

Face~5 produces the elliptic $r$-matrix; Face~7 produces the
commuting Hamiltonians it generates. At genus~$0$, Face~7 identifies
the collision residue with the Gaudin Hamiltonians
$H_i^{\mathrm{Gaudin}}$ of Theorem~\ref{thm:hdm-face-7}. At
genus~$1$, the Gaudin Hamiltonians become sums over
$\xi_\tau(z_i - z_j)$ weighted by the Casimir $\Omega_{ij}$, and
their commutativity follows from the elliptic classical Yang--Baxter
equation~\eqref{eq:g1sf-elliptic-cybe} proved in Face~5. The logical
chain is: propagator $\to$ collision residue $\to$ $r$-matrix $\to$
Yang--Baxter $\to$ commuting Hamiltonians. Each step is forced by the
preceding one.

\begin{theorem}[Face~7 at genus~$1$: elliptic Gaudin Hamiltonians;
\ClaimStatusConditional]
\label{thm:g1sf-face-7}
\index{Gaudin Hamiltonians!elliptic|textbf}
\index{seven faces!genus-1 Face 7|textbf}
\textbf{Type signature.} (Open quadrant; bar/twisting presentation on
$(E_\tau, D, \tau)$; level~$4$ for the Gaudin/Hamiltonian projection
of the level-$2$ collision residue; hypothesis package:
$\cA = \widehat{\fg}_k$ with $k\neq -h^\vee$,
KZ/Sugawara level prefix; the elliptic classical Yang--Baxter
equation~\eqref{eq:g1sf-elliptic-cybe}.)
For $\cA = \widehat{\fg}_k$, the elliptic Gaudin Hamiltonians on
$\mathrm{Conf}_n(E_\tau)$ are
\begin{equation}\label{eq:g1sf-elliptic-gaudin}
H_i^{\mathrm{ell}}
\;=\;
\sum_{j \neq i}
\bigl[r_{\widehat{\fg}_k}^{(1)}(z_i - z_j, \tau)\bigr]_{ij}
\;=\;
\frac{k}{k+h^\vee}
\sum_{j \neq i}
\bigl[r^{\mathrm{ell}}_\fg(z_i-z_j,\tau)\bigr]_{ij}.
\end{equation}
Its Cartan/KZB projection is
$\frac{k}{k+h^\vee}\sum_{j\neq i}\Omega_{ij}\xi_\tau(z_i-z_j)$.
They mutually commute:
$[H_i^{\mathrm{ell}}, H_j^{\mathrm{ell}}] = 0$ for all $i, j$.
The commutativity follows from the classical elliptic Yang--Baxter
equation~\eqref{eq:g1sf-elliptic-cybe}.
\end{theorem}

\begin{proof}
The identification of $H_i^{\mathrm{ell}}$ with the collision residue
is immediate from Theorem~\ref{thm:g1sf-face-1} in the
KZ/Sugawara-prefixed convention (Lemma~\ref{lem:g1sf-casimir-bridge}):
$r_{\widehat{\fg}_k}^{(1)}(z, \tau) = k\,r^{\mathrm{ell}}_\fg(z,\tau)/(k+h^\vee)$,
and the Gaiotto--Zeng formula~\eqref{eq:hdm-face-4} produces the
sum~\eqref{eq:g1sf-elliptic-gaudin}.

For commutativity, the modular MC equation at genus~$1$ gives
$D_\cA \Theta_\cA + \frac{1}{2}[\Theta_\cA, \Theta_\cA] = 0$
restricted to the collision-depth-$2$ stratum of $\overline{\cM}_{1,n}$.
Expanding the bracket term and projecting onto the $(i,j)$ and
$(i,k)$ channels produces the full elliptic infinitesimal braid
identity
\begin{equation}\label{eq:g1sf-elliptic-ibr}
[r^{\mathrm{ell}}_{ij}(z_{ij}),\;
r^{\mathrm{ell}}_{ik}(z_{ik})
+r^{\mathrm{ell}}_{jk}(z_{jk})]
\;=\;
0,
\end{equation}
whose Cartan/KZB projection is the $\xi_\tau\Omega$ channel. This
follows from the classical elliptic YBE~\eqref{eq:g1sf-elliptic-cybe}
by specializing to three points. The
commutativity $[H_i^{\mathrm{ell}}, H_j^{\mathrm{ell}}] = 0$ is the
$n$-point consequence of~\eqref{eq:g1sf-elliptic-ibr}, exactly as in
the genus-$0$ case (Theorem~\ref{thm:hdm-face-7}).
\end{proof}

\begin{remark}[Elliptic Gaudin and the Hitchin system]
\label{rem:g1sf-hitchin}
The elliptic Gaudin model is the genus-$1$ instance of the Hitchin
integrable system on $\mathrm{Bun}_G(E_\tau)$. The Hitchin
Hamiltonians are sections of a line bundle on the base of the Hitchin
fibration; at genus~$1$, the base is the affine space of elliptic
spectral curves, and the Hitchin Hamiltonians restrict to the
elliptic Gaudin Hamiltonians on the trivial $G$-bundle. The
genus-$1$ identification is the proved elliptic comparison surface.
For $g\geq2$, the prime-form collision residue gives the expected
Hitchin comparison kernel on $\mathrm{Bun}_G(\Sigma_g)$; the equality
with the full higher-genus Hitchin system is the conjectural content of
Section~\ref{sec:frontier-modular-holography-platonic}.
\end{remark}


%% ====================================================================
%% 7. CLASS M AT GENUS 1: HIGHER-ORDER ELLIPTIC OPERATORS
%% ====================================================================

\section{Class $M$ at genus~$1$: higher-order elliptic operators}
\label{sec:g1sf-class-m}

Faces~1, 4, 5, 7 all concern affine Kac--Moody algebras (class~$L$),
where the collision expansion terminates at
$c_0 = k\,\Omega/(k+h^\vee)$: collision depth $k_{\max} = 1$, and
only the $\xi_\tau$-term survives in the Cartan/KZB projection. The
classical affine elliptic $r$-matrix, KZB, and Gaudin lanes live at
this depth.

Class~$M$ algebras (Virasoro, $\cW_N$) break this barrier. Their
OPEs have poles of order $\geq 4$, and the collision expansion
involves $P_\tau$, $P_\tau'$, and higher derivatives of the
prime-form curvature kernel. The resulting genus-$1$ collision
residues are elliptic differential operators of order $\geq 2$, forced
by OPE collision depth beyond the affine Casimir $r$-matrix lane.

\subsection{Virasoro at genus~$1$}
\label{sec:g1sf-virasoro}

The Virasoro algebra has $p_{\max} = 4$ (the $T(z)\, T(w)$ OPE has a
pole of order~$4$). the bar propagator absorbs one order,
giving collision depth $k_{\max} = 3$. The genus-$0$ collision
residue is
\begin{equation}\label{eq:g1sf-vir-genus0}
r_{\mathrm{Vir}_c}(z)
\;=\;
\frac{c/2}{z^3}
\;+\;
\frac{2T}{z},
\end{equation}
with $c_0 = 2T$ (the stress-energy tensor acting on conformal
blocks), $c_1 = 0$ (the $z^{-2}$ pole is absent from the $r$-matrix
by the parity argument: the $d\log$ extraction sends
$z^{-2n}$ to $z^{-(2n-1)}$), and $c_2 = c/2$.

\begin{theorem}[Virasoro genus-$1$ collision residue;
\ClaimStatusConditional]
\label{thm:g1sf-virasoro}
\index{Virasoro!genus-1 collision residue|textbf}
\index{collision residue!Virasoro at genus 1|textbf}
\textbf{Type signature.} (Open quadrant; bar/twisting presentation on
$(E_\tau, D, \tau)$; level~$2$ for the elliptic collision residue,
level~$4$ for its action on conformal blocks via Ward identity;
hypothesis package: Virasoro at central charge~$c$, completed
bar coalgebra, prime-form normalisation. The $c=13$ KSDual sublocus
is the $\Z/2$-fixed point under
$\mathrm{Vir}_c\mapsto\mathrm{Vir}_{26-c}$
\textup{(}Remark~\textup{\ref{rem:c13-ksdual-poincare}}\textup{)}.)
The genus-$1$ collision residue of $\mathrm{Vir}_c$ is:
\begin{equation}\label{eq:g1sf-vir-genus1}
r_{\mathrm{Vir}_c}^{(1)}(z, \tau)
\;=\;
2T \cdot \xi_\tau(z)
\;-\;
\frac{c}{4}\, P_\tau'(z).
\end{equation}
This is an elliptic differential operator of order~$2$ acting on
conformal blocks on $E_\tau$: the $\xi_\tau$-term acts by
multiplication, while the $P_\tau'$-term acts as a second-order
differential operator through the stress-energy Ward identity.
\end{theorem}

\begin{proof}
The genus-$0$ collision expansion has $c_0 = 2T$, $c_1 = 0$, and
$c_2 = c/2$. By
Theorem~\ref{thm:g1sf-elliptic-regularization}, the elliptic
regularization~\eqref{eq:g1sf-elliptic-expansion} gives
\[
r_{\mathrm{Vir}_c}^{(1)}(z, \tau)
\;=\;
c_0\, \xi_\tau(z)
+ c_1\, P_\tau(z)
- \tfrac{1}{2}\, c_2\, P_\tau'(z).
\]
Since $c_1 = 0$, the $P_\tau$-term vanishes, and the result
is~\eqref{eq:g1sf-vir-genus1}. The Laurent expansion near
$z = 0$ gives
$2T/z - (c/4)\cdot(-2/z^3 + O(z)) = (c/2)/z^3 + 2T/z + \cdots$;
the leading poles $(c/2)/z^3$ and $2T/z$ match the genus-$0$ collision
residue~\eqref{eq:g1sf-vir-genus0} at orders $z^{-3}$ and $z^{-1}$
respectively, consistent with $c_2 = c/2$ and the Laurent expansion
$\wp'(z) = -2/z^3 + O(z)$.

The $P_\tau'$-term acts as a second-order differential operator on
conformal blocks because the Virasoro OPE at order $(z-w)^{-4}$
involves $c/2 \cdot \mathrm{id}$, which through the Ward identity
produces a second derivative $\partial_z^2$ acting on the
correlator. The elliptic regularization replaces $1/z^4$ by
$-P_\tau'(z)/2$, and the result is a second-order
elliptic differential operator.
\end{proof}

\begin{remark}[The $c_1 = 0$ vanishing]
\label{rem:g1sf-c1-vanishing}
The vanishing of $c_1$ in the Virasoro collision expansion is the
parity phenomenon: the OPE mode $T_{(2)}T = \partial T$ is a
total derivative, and its contribution to the $r$-matrix vanishes
after the $d\log$-extraction. In the elliptic regularization, this
means the $P_\tau(z)$-term is absent from the Virasoro genus-$1$
collision residue. The first non-trivial correction beyond
$\xi_\tau$ involves $P_\tau'=\wp'$, not $P_\tau$.
\end{remark}

\subsection{$\cW_N$ at genus~$1$}
\label{sec:g1sf-wn}

The $\cW_N$ algebra has generators $W_2 = T, W_3, \ldots, W_N$ of
conformal weights $2, 3, \ldots, N$. The maximal OPE pole order among
generating fields is $p_{\max} = 2N$ (from the $W_N$-$W_N$ OPE), giving
collision depth $k_{\max} = 2N - 1$.

\begin{theorem}[$\cW_N$ genus-$1$ collision residue; \ClaimStatusConditional]
\label{thm:g1sf-wn}
\index{W-algebra@$\cW_N$!genus-1 collision residue|textbf}
\textbf{Type signature.} (Open quadrant; bar/twisting presentation on
$(E_\tau, D, \tau)$; level~$2$ for the elliptic collision residue,
level~$4$ for the higher-order elliptic differential operators it
generates; hypothesis package: principal $\cW_N$ at central charge $c$,
weight-completed bar coalgebra, prime-form normalisation,
finite collision depth $k_{\max} = 2N{-}1$.)
The genus-$1$ collision residue of $\cW_N$ involves the prime-form
curvature kernel up to $P_\tau^{(2N-3)}$:
\begin{equation}\label{eq:g1sf-wn-genus1}
r_{\cW_N}^{(1)}(z, \tau)
\;=\;
c_0\, \xi_\tau(z)
\;+\;
\sum_{\substack{n = 1 \\ c_n \neq 0}}^{2N-2}
\frac{(-1)^{n-1}}{n!}\,
c_n \cdot P_\tau^{(n-1)}(z),
\end{equation}
where $c_n$ are the OPE-mode collision coefficients of $\cW_N$. The
highest-order term is an elliptic differential operator of order
$2N - 2$. For $\cW_3$: the collision expansion involves
$\xi_\tau, P_\tau, P_\tau', P_\tau'', P_\tau'''$, producing fourth-order
elliptic differential operators.
\end{theorem}

\begin{proof}
Direct specialization of
Theorem~\ref{thm:g1sf-elliptic-regularization} to the $\cW_N$ OPE
data with $p_{\max} = 2N$. The collision depth $k_{\max} = 2N - 1$
means the genus-$0$ expansion
$r_{\cW_N}(z) = \sum_{n=0}^{2N-2} c_n/z^{n+1}$ has $2N - 1$ terms.
Each term is regularized by the corresponding elliptic function:
$1/z \mapsto \xi_\tau(z)$ and
$1/z^{n+1} \mapsto (-1)^{n-1}P_\tau^{(n-1)}(z)/n!$ for $n \geq 1$.
\end{proof}

\begin{remark}[Higher-order elliptic operators]
\label{rem:g1sf-higher-order-operators}
\index{collision residue!higher-order elliptic operators|textbf}
The elliptic differential operators produced by the class~$M$
collision residue at genus~$1$ begin where the affine Casimir lane
stops. The Hitchin and KZB comparison kernels use the depth-$1$
prime-form term and the depth-$2$ curvature kernel. Class~$M$
collision residues at depth~$\geq3$ involve
$P_\tau',P_\tau'',\ldots,P_\tau^{(2N-3)}$; these terms are the elliptic
images of the higher OPE poles of Virasoro and $\cW$-algebras.
\end{remark}

\begin{remark}[Class $G$ and class $C$ at genus~$1$]
\label{rem:g1sf-class-gc}
For class~$G$ (Heisenberg, lattice VOAs) with
$p_{\max} = 2$, $k_{\max} = 1$, the genus-$1$ collision residue
has the same form as class~$L$: a single $\xi_\tau$-term. For
class~$C$ ($\beta\gamma$) with $p_{\max} = 1$, $k_{\max} = 0$, the
genus-$0$ collision residue vanishes, and the genus-$1$ collision
residue is identically zero: $r_{\beta\gamma}^{(1)}(z, \tau) = 0$.
The binary collision residue of the $\beta\gamma$ system is zero. The
shadow class hierarchy
$G, L, C, M$ is unchanged at genus~$1$; only the concrete
realization changes from rational to elliptic functions.
\end{remark}


%% ====================================================================
%% 8. THE GENUS-1 SEVEN-FACE MASTER THEOREM
%% ====================================================================

\section{The genus-$1$ seven-face master theorem}
\label{sec:g1sf-master}

The seven faces at genus~$0$ are seven extraction maps to one Laurent
kernel. At genus~$1$, they are seven extraction maps to one
prime-form-normalized elliptic kernel. Faces~1, 4, 5, 7 are proved on
the affine Kac--Moody comparison surface above. Faces~2, 3, 6 enter
when their elliptic comparison functors are fixed; the theorem states
the common extracted kernel shared by the source categories, brackets,
connections, and Hamiltonian algebras after their extraction maps are
applied.

\begin{theorem}[Genus-$1$ seven-face identification for affine Kac--Moody;
\ClaimStatusConditional]
\label{thm:g1sf-master}
\index{seven faces!genus-1 master theorem|textbf}
\index{collision residue!genus-1 seven-face identification|textbf}
\textbf{Type signature.} (Open quadrant; spans levels~$2$
\textup{(}twisting morphism\textup{)}, $4$ \textup{(}Hamiltonians and
brackets\textup{)}, and~$5$ \textup{(}modular shadow $d\tau$-direction
and $\mathrm{SL}(2,\Z)$-equivariance\textup{)}; hypothesis package:
$\cA=\widehat{\fg}_k$, $k\neq -h^\vee$, KZ/Sugawara level prefix per
Lemma~\ref{lem:g1sf-casimir-bridge}, comparison maps for Faces F2, F3,
F6 fixed when those are invoked.)

For $\cA = \widehat{\fg}_k$ at level $k \neq -h^\vee$, let
$I^{(1)}(\cA)$ be the set of genus-$1$ faces whose comparison
hypotheses hold. For each $i\in I^{(1)}(\cA)$ let
$\mathsf E_i^{(1)}$ denote the face-specific extraction map to the
binary elliptic collision kernel. Then
\[
\mathsf E_i^{(1)}(F_i^{(1)})
\;=\;
r_\cA^{(1)}(z,\tau)
\qquad (i\in I^{(1)}(\cA)),
\]
where
$r_\cA^{(1)}(z, \tau)=
\operatorname{Res}^{\mathrm{coll}}_{1,2}(\Theta_\cA)|_{E_\tau}$.
The seven affine genus-$1$ faces are:
\begin{enumerate}[label=\textup{(F\arabic*${}^{(1)}$)}]
\item \textbf{Elliptic twisting morphism.}
$r_\cA^{(1)} = \pi_\cA^{(1)}$, the genus-$1$ bar-cobar twisting
morphism \textup{(}Theorem~\textup{\ref{thm:g1sf-face-1})}.

\item \textbf{Elliptic line-operator $r$-matrix.}
$\mathsf E_2^{(1)}(r^{\mathrm{DNP},(1)})=r_\cA^{(1)}$ for the
elliptic spectral kernel of the line-operator category on
$E_\tau \times \mathbb{R}$ in the DNP framework~\cite{DNP25}.

\item \textbf{Elliptic $\lambda$-bracket.}
$\mathsf E_3^{(1)}$ sends the elliptic $\lambda$-bracket of the
classical PVA on $E_\tau$ to the kernel $r_\cA^{(1)}$
\textup{(}Khan--Zeng~\cite{KhanZeng25}, genus-$1$ specialization\textup{)}.

\item \textbf{KZB connection.}
The $dz$-component of the KZB connection equals the collision residue;
the $d\tau$-component is its $z$-derivative
\textup{(}Theorem~\textup{\ref{thm:g1sf-face-4})}.

\item \textbf{Belavin--Drinfeld elliptic $r$-matrix.}
$r_\cA^{(1)} = k\,r^{\mathrm{ell}}_\fg/(k + h^\vee)$
\textup{(}Theorem~\textup{\ref{thm:g1sf-face-5})}; the level
prefix~$k$ is required.

\item \textbf{Elliptic Sklyanin bracket.}
The classical-limit extraction from the elliptic Sklyanin Poisson
bracket on $\fg^*$ gives the Belavin--Drinfeld kernel, generalizing
the genus-$0$ Sklyanin bracket of
Semenov-Tian-Shansky~\cite{STS83}.

\item \textbf{Elliptic Gaudin Hamiltonians.}
$H_i^{\mathrm{ell}} = \sum_{j \neq i}\,[r_\cA^{(1)}(z_i - z_j)]_{ij}$
\textup{(}Theorem~\textup{\ref{thm:g1sf-face-7})}.
\end{enumerate}
All available extracted kernels agree as elements of
$\End_\cA(2)[\![z^{-1}]\!] \otimes \mathbb{C}[\![\tau]\!]$:
\begin{equation}\label{eq:g1sf-master-equation}
\boxed{\;
r_\cA^{(1)}(z, \tau)
\;=\;
\mathsf E_1^{(1)}(\pi_\cA^{(1)})
\;=\;
\mathsf E_2^{(1)}(r^{\mathrm{DNP},(1)})
\;=\;
\mathsf E_3^{(1)}\!\left(
\mathrm{PVA}^{\mathrm{ell}}_{\lambda}(\cA)\right)
\;=\;
\mathsf E_4^{(1)}(H_z^{\mathrm{KZB}})
\;=\;
\frac{k\,r^{\mathrm{ell}}_\fg}{k + h^\vee}
\;=\;
\mathsf E_6^{(1)}(\{\cdot,\cdot\}_{\mathrm{Sk}}^{\mathrm{ell}})
\;=\;
\mathsf E_7^{(1)}(H_i^{\mathrm{ell}})
\;}
\end{equation}
with
\[
r_\cA^{(1)}(z,\tau)
=
c_0\xi_\tau(z)
+
\sum_{n\geq1}\frac{(-1)^{n-1}}{n!}c_nP_\tau^{(n-1)}(z)
\]
in the prime-form normalization of
Theorem~\ref{thm:g1sf-elliptic-regularization}.
\end{theorem}

\begin{proof}
Face~F1$^{(1)}$ is Theorem~\ref{thm:g1sf-face-1}.
Face~F4$^{(1)}$ is Theorem~\ref{thm:g1sf-face-4}.
Face~F5$^{(1)}$ is Theorem~\ref{thm:g1sf-face-5}.
Face~F7$^{(1)}$ is Theorem~\ref{thm:g1sf-face-7}.
Faces~F2$^{(1)}$, F3$^{(1)}$, F6$^{(1)}$ use the same comparison maps
as the genus-$0$ identification theorems
(Theorems~\ref{thm:hdm-face-2},~\ref{thm:hdm-face-3},~\ref{thm:hdm-face-6}),
with the rational kernel replaced by the prime-form kernel
$\xi_\tau$ and its curvature derivatives $P_\tau^{(m)}$. The
F3$^{(1)}$ Khan--Zeng PVA comparison is used only through this
binary-kernel extraction. The transitive closure of the extracted
kernel equalities yields~\eqref{eq:g1sf-master-equation}.
\end{proof}

\begin{remark}[Scope: class~$M$]
\label{rem:g1sf-class-m-scope}
For class~$M$ algebras (Virasoro, $\cW_N$), Faces~F1$^{(1)}$
and~F4$^{(1)}$ are defined on the modular-Koszul locus by the
bar-cobar twisting morphism and the level-$5$ modular shadow
connection. Faces~F5$^{(1)}$--F7$^{(1)}$ are affine Kac--Moody
comparison faces: the Virasoro and $\cW_N$ residues of
Theorems~\ref{thm:g1sf-virasoro} and~\ref{thm:g1sf-wn} carry
higher-order $P_\tau^{(m)}$ terms that are not classical Casimir
$r$-matrices. Their commutativity, where defined, is a consequence of
the genus-$1$ MC equation; the modular ($\mathrm{SL}(2,\Z)$) shadow on
the $d\tau$-direction is the level-$5$ trace + clutching package on
the open factorization category, not an intrinsic closed-algebra
modularity. On the \emph{KSDual} sublocus
($\Z/2$-fixed under $\cA\mapsto\cA^!$), Theorem~A becomes an
equivalence and the genus-$1$ shadow strengthens accordingly: for
Virasoro this sublocus is $c=13$
(cf.~Remark~\ref{rem:c13-ksdual-poincare} in
Chapter~\ref{chap:NAP-computations}); for affine Kac--Moody it is the
Feigin--Frenkel level $-2k-2h^\vee = k$ specialisation, which is
empty (no real-level fixed point), reflecting the fact that the
generic affine Kac--Moody class~$L$ row is not a self-dual KSDual
locus and the modular reconstruction at genus~$1$ retains its
adjunction-with-comparison-map form.
\end{remark}


%% ====================================================================
%% 9. DEGENERATION TO GENUS 0
%% ====================================================================

\section{Degeneration to genus~$0$}
\label{sec:g1sf-degeneration}

The genus-$1$ identification must recover the genus-$0$ master theorem
(Theorem~\ref{thm:hdm-seven-way-master}) in the degeneration limit
$\tau \to i\infty$. In this limit the elliptic curve $E_\tau$
pinches to a nodal rational curve: the period $\tau$ becomes infinite,
the lattice $\mathbb{Z} + \mathbb{Z}\tau$ degenerates, and every
elliptic function reduces to a rational one. The recovery is uniform
across all seven faces.

\begin{theorem}[Degeneration to genus~$0$; \ClaimStatusConditional]
\label{thm:g1sf-degeneration}
\index{collision residue!degeneration to genus 0}
\textbf{Type signature.} (Open quadrant; bar/twisting presentation on
$(E_\tau, D, \tau)$; level~$2$ degeneration to genus~$0$ collision
residue; hypothesis package: $q = e^{2\pi i\tau}\to 0$ limit on the
$k\neq 0,-h^\vee$ normalization locus per
Lemma~\ref{lem:g1sf-casimir-bridge}.)
As $\operatorname{Im}(\tau) \to \infty$
\textup{(}equivalently, $q = e^{2\pi i\tau} \to 0$\textup{)}, all
seven genus-$1$ faces degenerate to their genus-$0$ counterparts:
\begin{enumerate}[label=\textup{(\roman*)}]
\item \textup{(}Propagator.\textup{)}
$\xi_\tau(z) \to \pi\cot(\pi z) \to 1/z + O(z)$.

\item \textup{(}Prime-form curvature kernels.\textup{)}
$P_\tau(z) \to \pi^2/\sin^2(\pi z) \to 1/z^2 + O(1)$, and
$P_\tau^{(m)}(z) \to (-1)^m(m+1)!/z^{m+2} + O(z^{-m})$ at leading
polar order.

\item \textup{(}KZB $\to$ KZ.\textup{)}
The $dz$-component of the level-prefixed KZB kernel degenerates to the
level-prefixed KZ kernel: $\xi_\tau(z_{ij}) \to 1/(z_{ij})$. The
$d\tau$-component vanishes because the modulus $\tau$ is not a
parameter at genus~$0$.

\item \textup{(}Elliptic $r$-matrix $\to$ rational $r$-matrix.\textup{)}
The Belavin $r$-matrix degenerates to the classical rational Casimir
kernel $r^{\mathrm{rat}}_\fg(z)$, through the
intermediate trigonometric $r$-matrix:
$r^{\mathrm{ell}} \to r^{\mathrm{trig}} \to r^{\mathrm{rat}}$.
The collision residue~\eqref{eq:g1sf-elliptic-rmatrix} inherits
the KZ/Sugawara-prefixed level normalization
\textup{(}Lemma~\ref{lem:g1sf-casimir-bridge}\textup{)}:
$r^{(1)}_{\widehat{\fg}_k}(z, \tau) \to k\,\Omega_{\mathrm{KZ}}/((k+h^\vee)z)$.
% Level prefix k/(k+h^v) in the KZ/Sugawara-prefixed convention; at k=0 the
% prefixed r-matrix vanishes (distinct from the unprefixed KZ-convention
% residue). See Lemma g1sf-casimir-bridge for reconciliation.
At the critical level $k = -h^\vee$ the normalized rational
$r$-matrix is singular \textup{(}Feigin--Frenkel centre
enlarges\textup{)}; at $k = 0$ the KZ/Sugawara-prefixed collision
residue $r^{(1)}_{\widehat{\fg}_0}$ vanishes, reflecting the
Drinfeld-normalization vanishing $r^{\mathrm{Dr}}_0 = 0$, not to be
confused with the KZ-normalization object
$\Omega_{\mathrm{KZ}}/(h^\vee z)$ of Theorem~\ref{thm:hdm-face-5},
which is a distinct spectral $r$-matrix
\textup{(}Lemma~\ref{lem:g1sf-casimir-bridge}(a)\textup{)}.
\item \textup{(}Elliptic Gaudin $\to$ rational Gaudin.\textup{)}
$H_i^{\mathrm{ell}} \to H_i^{\mathrm{Gaudin}}$, the standard rescaled
Gaudin Hamiltonian of Theorem~\ref{thm:hdm-face-7}.
\end{enumerate}
In particular, the genus-$1$ master
equation~\eqref{eq:g1sf-master-equation} degenerates to the genus-$0$
master equation~\eqref{eq:hdm-master-equation} in the limit
$q \to 0$.
\end{theorem}

\begin{proof}
As $\operatorname{Im}(\tau) \to \infty$, $q = e^{2\pi i\tau} \to 0$,
and the prime-form kernels have the standard degenerations. The
logarithmic kernel: $\xi_\tau(z) = \pi\cot(\pi z) + O(q)$, and
$\pi\cot(\pi z) = 1/z - 2z\sum_{n=1}^\infty 1/(z^2 - n^2)$, which
gives $1/z + O(z)$ for $|z| < 1$. The curvature kernel:
$P_\tau(z) = \pi^2/\sin^2(\pi z) + O(q)$,
and $\pi^2/\sin^2(\pi z) = 1/z^2 + O(1)$. Higher derivatives
degenerate analogously.

The KZB degeneration: with $\xi_\tau(z) \to 1/z$, the
$dz$-component~\eqref{eq:g1sf-kzb-z-component} becomes
$\sum_{j \neq i}k\,\Omega_{ij}/((k+h^\vee)(z_i - z_j))\, dz_{ij}$,
which is the level-prefixed KZ kernel. The unprefixed KZ convention is
recovered on the $k\neq0$ normalization locus as in
Lemma~\ref{lem:g1sf-casimir-bridge}. The $d\tau$-component has no
genus-$0$ counterpart.

The elliptic $r$-matrix degeneration: as $q \to 0$,
$\phi_\pm(z, \tau) \to \pi/\sin(\pi z) + O(q)$, recovering the
trigonometric $r$-matrix. A further limit $z \mapsto z/L$,
$L \to \infty$ gives $\pi/\sin(\pi z/L) \to L/z$, recovering the
rational $r$-matrix
$k\,\Omega_{\mathrm{KZ}}/((k+h^\vee)z) = k\,r^{\mathrm{rat}}_\fg(z)/(k+h^\vee)$
in the KZ/Sugawara-prefixed normalization used throughout this chapter
% Level prefix k/(k+h^v); coupling-rescaled from Drinfeld k*Omega_tr/z
% through Lemma g1sf-casimir-bridge, and distinct at k=0.
\textup{(}cf.\ Eq.~\eqref{eq:g1sf-elliptic-rmatrix} and the
KZ/Sugawara-prefixed presentation of the genus-$0$ collision residue
$r_{\widehat{\fg}_k}(z) = k\,\Omega_{\mathrm{KZ}}/((k+h^\vee)z)$,
which is the explicit-level-prefix variant of the KZ-convention
residue $\Omega_{\mathrm{KZ}}/((k+h^\vee)z)$ of
Theorem~\ref{thm:hdm-face-5}; the two are not equal as functions of
$k$ and differ decisively at $k = 0$\textup{)}.
Both limits are continuous in the topology of formal Laurent series
for $k \neq 0, -h^\vee$.
\end{proof}

The degeneration theorem identifies the genus-$0$ Laurent kernel as
the polar limit of the genus-$1$ prime-form kernel. At $q=0$, the
elliptic curve degenerates to the nodal rational curve, and each
genus-$1$ extraction map reduces to its genus-$0$ extraction map.


%% ====================================================================
%% 10. COMPUTATIONAL VERIFICATION
%% ====================================================================

\section{Computational verification}
\label{sec:g1sf-compute}

The genus-$1$ seven-face identifications are verified by the compute
engine \texttt{theorem\_genus1\_seven\_faces\_engine.py}, which
implements the following tests.

\begin{enumerate}[label=(\arabic*)]
\item \textbf{Elliptic regularization.}
For each standard family $\cA$ in the landscape census
(Chapter~\ref{ch:landscape-census}), verify that the genus-$1$
collision residue~\eqref{eq:g1sf-elliptic-expansion} matches the
elliptic regularization of the genus-$0$ collision
expansion at $O(q^0)$ to machine precision.

\item \textbf{KZB identification.}
For $\cA = \widehat{\mathfrak{sl}}_2$ at levels
$k = 1, 2, 3, 5, 10$, verify that the $dz$-component of the
KZB connection equals the collision residue, and that the
$d\tau$-component equals $-\partial_z r_\cA^{(1)}$ in the
prime-form normalization.

\item \textbf{Elliptic CYBE.}
Verify the classical elliptic Yang--Baxter
equation~\eqref{eq:g1sf-elliptic-cybe} numerically for
$\mathfrak{sl}_2$ at $20$ random evaluation points
$(z_1, z_2, z_3) \in E_\tau^3$ and $5$ values of $\tau$.

\item \textbf{Gaudin commutativity.}
Verify $[H_i^{\mathrm{ell}}, H_j^{\mathrm{ell}}] = 0$ for
$\mathfrak{sl}_2$ on $\mathrm{Conf}_3(E_\tau)$ at $10$ random
configurations and $5$ values of $\tau$.

\item \textbf{Degeneration.}
Verify that
$\|r_\cA^{(1)}(z, \tau) - r_\cA(z)\| / \|r_\cA(z)\| < 10^{-6}$
for $\operatorname{Im}(\tau) \geq 12$ across all standard
families.

\item \textbf{Class~$M$ structure.}
For $\mathrm{Vir}_c$ at $c = 1, 1/2, 25, 26$, verify the
collision residue~\eqref{eq:g1sf-vir-genus1} and confirm that
the $c_1 = 0$ vanishing holds. For $\cW_3$ at
$c = 2, 50, 98$, verify the collision
residue~\eqref{eq:g1sf-wn-genus1} and confirm the appearance
of terms through $P_\tau'''$.
\end{enumerate}


%% ====================================================================
%% 11. THE B-CYCLE MONODROMY AND THE QUANTUM GROUP
%% ====================================================================

\section{The $B$-cycle monodromy and the quantum group parameter}
\label{sec:g1sf-b-cycle}

The genus-$1$ collision residue carries information that is invisible
at genus~$0$: the $B$-cycle monodromy. For affine Kac--Moody, the
$B$-cycle monodromy of the KZB connection produces the quantum group
half-monodromy parameter $q_{\mathrm{QG}} = \exp(\pi i/(k + h^\vee))$
(Theorem~\ref{thm:b-cycle-quantum-group}).

\begin{theorem}[$B$-cycle monodromy of the collision residue;
\ClaimStatusProvedHere]
\label{thm:g1sf-b-cycle-monodromy}
\index{B-cycle monodromy@$B$-cycle monodromy!collision residue|textbf}
\textbf{Type signature.} (Open quadrant; level~$5$ for the
$\mathrm{SL}(2,\Z)$-equivariant $B$-cycle datum on
$\overline{\cM}_{1,2}$; hypothesis package: $\cA$ modular Koszul
with $c_0\neq 0$; for the affine specialization,
$\cA=\widehat{\fg}_k$ with $k\neq -h^\vee$ so that $c_0=k\Omega/(k+h^\vee)$
is finite. The quantum-group parameter
$q_{\mathrm{QG}} = \exp(\pi i/(k+h^\vee))$ is the level-$5$ modular
reconstruction of the half-monodromy on the open category, not an
intrinsic closed-algebra invariant.)

Let $\cA$ be a modular Koszul chiral algebra with $c_0 \neq 0$. The
$B$-cycle monodromy of the genus-$1$ collision residue is
\begin{equation}\label{eq:g1sf-b-cycle}
r_\cA^{(1)}(z + \tau, \tau) - r_\cA^{(1)}(z, \tau)
\;=\;
-2\pi i \cdot c_0.
\end{equation}
For $\cA = \widehat{\fg}_k$, $c_0 = k\,\Omega/(k + h^\vee)$, and the
prime-form monodromy is $-2\pi i\, k\,\Omega/(k + h^\vee)$. The
full exponentiated $B$-cycle monodromy of the KZB connection is
$q_{\mathrm{QG}}^{2H}$; the standard quantum group parameter is the
half-monodromy:
\begin{equation}\label{eq:g1sf-quantum-parameter}
q_{\mathrm{QG}}
\;=\;
\exp\!\Bigl(\frac{\pi i}{k + h^\vee}\Bigr),
\end{equation}
recovering the standard relation between the level and the quantum
group deformation parameter.
\end{theorem}

\begin{proof}
The $B$-cycle monodromy of $r_\cA^{(1)}$ is computed
in~\eqref{eq:g1sf-B-cycle-monodromy}: it equals
$-2\pi i \cdot c_0$, and the doubly-periodic terms do not
contribute. For $\cA = \widehat{\fg}_k$, the full exponentiated
monodromy of the KZB connection around the $B$-cycle $\beta_i$
is $\exp(2\pi i \hbar H_i)=q_{\mathrm{QG}}^{2H_i}$ where
$\hbar = 1/(k + h^\vee)$ is the Drinfeld parameter and $H_i$ acts
by the weight of the $i$-th representation. Thus
$q_{\mathrm{QG}} = \exp(\pi i\hbar) = \exp(\pi i/(k + h^\vee))$. The
computation is the content of
Theorem~\ref{thm:b-cycle-quantum-group}(ii).
\end{proof}

\begin{remark}[$B$-cycle monodromy and spectral data]
\label{rem:g1sf-b-cycle-face-0}
The $B$-cycle monodromy is a genus-$1$ phenomenon: it
has no genus-$0$ counterpart. The quantum group parameter
$q_{\mathrm{QG}}$ is extracted from the same prime-form collision
residue that produces the KZB connection, the elliptic $r$-matrix, and
the elliptic Gaudin Hamiltonians. The MC element $\Theta_\cA$ on
$E_\tau$ encodes both the spectral data and the $B$-cycle monodromy.
\end{remark}

\begin{remark}[Class~$M$ at genus~$1$: Virasoro monodromy]
\label{rem:g1sf-class-m-no-qg}
For class~$M$ algebras, the $B$-cycle monodromy
$-2\pi i \cdot c_0$ is non-zero (since $c_0 = 2T \neq 0$ for
Virasoro). However, $c_0$ acts on conformal blocks by the zero mode
$L_0$. The Drinfeld quantum-group construction requires a
finite-dimensional Cartan action, so the exponentiated monodromy
$\exp(2\pi i \cdot L_0)$ lies in the level-$5$ modular shadow
connection on the open factorization category rather than in a
finite-dimensional closed quantum group. This is consistent with the
master critique: closed shadows of class-$M$ algebras carry modular
consequences (Virasoro $c\mapsto 26-c$ Verdier complementarity, Zhu
algebra rationality on the rational locus), but those modular
consequences are reconstructed from the open trace + clutching, not
from an intrinsic closed-algebra modularity claim.
\end{remark}


%% ====================================================================
%% 12. HIGHER-GENUS PRIME-FORM KERNELS
%% ====================================================================

\section{Higher-genus prime-form kernels}
\label{sec:g1sf-outlook}

At genus $g \geq 2$, the holomorphic polar kernel is
$d\log E(z,w;\Sigma_g)$, the logarithmic derivative of Fay's prime
form on $\Sigma_g$. The bar-intrinsic MC element still restricts to a
collision residue
$r_\cA^{(g)}(z,w;\Sigma_g)=
\operatorname{Res}^{\mathrm{coll}}_{g,2}(\Theta_\cA)$.
Face~1 follows from the same bar-cobar twisting construction. The
Hitchin, spectral-$r$, Sklyanin, and Gaudin comparison faces require
additional higher-genus comparison maps.

\begin{conjecture}[Higher-genus seven-face identification]
\label{conj:g1sf-higher-genus}
\ClaimStatusConjectured
\index{seven faces!higher-genus|textbf}
\textbf{Type signature.} (Open quadrant; bar/twisting presentation on
$(\Sigma_g, D, \tau)$ for the prime-form collision residue; spans
levels~$2$ \textup{(}twisting morphism\textup{)}, $4$
\textup{(}brackets, Hamiltonians, $r$-matrix\textup{)}, and~$5$
\textup{(}modular shadow on $\overline{\cM}_{g,n}$\textup{)};
hypothesis package: modular Koszul $\cA$, prime-form collision
kernel $d\log E(z,w;\Sigma_g)$ on $\Sigma_g$, higher-genus comparison
maps for Faces~F2--F7.)
For every modular Koszul chiral algebra $\cA$ and every genus
$g \geq 0$, the collision residue
$r_\cA^{(g)} = \operatorname{Res}^{\mathrm{coll}}_{g,2}(\Theta_\cA)$
admits seven equivalent realizations that specialize, for
$\cA = \widehat{\fg}_k$, to: \textup{(1)} the genus-$g$ twisting
morphism, \textup{(2)} the line-operator $r$-matrix on
$\Sigma_g \times \mathbb{R}$, \textup{(3)} the classical
$\lambda$-bracket on $\Sigma_g$, \textup{(4)} the Hitchin
Hamiltonians on $\mathrm{Bun}_G(\Sigma_g)$, \textup{(5)} the spectral
$r$-matrix on $\Sigma_g$, \textup{(6)} the Sklyanin bracket on
$\Sigma_g$, and \textup{(7)} the Gaudin model on $\Sigma_g$.
\end{conjecture}

The proved cases are genus~$0$
(Theorem~\ref{thm:hdm-seven-way-master}) and genus~$1$
(Theorem~\ref{thm:g1sf-master}). The case $g\geq2$ is conjectural
outside Face~1 and the existence of the prime-form collision kernel.

\section{Structural anchors}
\label{sec:g1sf-anchors}

\begin{remark}[Structural anchors for the genus-$1$ seven faces]
\label{rem:g1sf-anchors}
\index{structural anchors!genus-1 seven-faces entry point}
The genus-$1$ seven-face theorem
(Theorem~\ref{thm:g1sf-master}) sits at the elliptic specialisation
of three structural anchors:
\begin{enumerate}[label=\textup{(W\arabic*)}]
\item \emph{Climax identity, elliptic.} The KZB connection of
 Face~F4 is the elliptic instance of the climax identity
 $d_{\mathrm{bar}} = \mathrm{KZ}^*(\nabla_{\mathrm{Arnold}})$
 of Chapter~\ref{ch:climax-theorem}: pulling back the elliptic
 Arnold connection
 $d-\xi_\tau(z_i-z_j)\Omega_{ij}\,d(z_i-z_j)$ through the KZB map
 produces the genus-$1$ bar differential. The Arakelov correction is
 the separate non-holomorphic term needed for single-valuedness on
 $E_\tau$; it is not part of the holomorphic polar residue.
\item \emph{Universal conductor.} The modular characteristic
 $\kappa(\widehat{\fg}_k) = \dim(\fg)(k+h^\vee)/(2h^\vee)$ enters
 the elliptic propagator's residue exactly through the level-shifted
 Sugawara denominator $1/(k+h^\vee)$, and the conductor identity
 $\kappa(\cA) = -c_{\mathrm{ghost}}(\BRST(\cA))$ of
 Chapter~\ref{chap:universal-conductor} provides the Lagrangian
 origin: the elliptic ghost system carries
 $c_{\mathrm{ghost}}^{\mathrm{WZW}} = -2\dim(\fg)$, and the
 cancellation with the matter conformal anomaly is the Sugawara
 nilpotence condition.
\item \emph{Koszul reflection at genus~$1$.} The genus-$1$
 bar-cobar adjunction
 (Theorem~B at $g=1$, with the Heisenberg case unconditional)
 implements the Koszul reflection theorem
 (Theorem~\ref{thm:koszul-reflection}, Vol~I Platonic Theorem~A)
 in the elliptic setting: the elliptic twisting morphism of
 Face~F1 is the unit of the genus-$1$ bar-cobar adjunction, and
 the seven elliptic faces are seven names for that unit datum
 after passage to the appropriate target category.
\end{enumerate}
The volume bridges run as follows. Vol~II's universal
celestial holography (Vol~II Theorem~\ref{V2-thm:universal-holography-master})
realises the elliptic seven faces on a 3d HT theory whose closed
colour is~$\widehat{\fg}_k$ and whose worldsheet is
$\bR_{\geq 0} \times E_\tau$; the elliptic line operators of Face~F2
are the holographic line defects in this 3d theory (level~$3$ vertical
equivalence to the $\mathsf{SC}^{\mathrm{ch,top}}$-brace bulk; level~$4$
identification with brane operators). Vol~III's
two-stage CY-to-chiral factorisation
$\Phi_3^{(\Sigma_2, E_\tau)} = \SpCh_{\Sigma_2, E_\tau} \circ \PhiFA_3$
realises the elliptic seven faces for
$\Phi_3^{(\Sigma_2,E_\tau)}(\mathrm{K3} \times E)$ in the BPS-counting
category, with Stage~1 (level~$1$) delivering the canonical
$E_3$-holomorphic factorisation algebra on $\mathrm{K3} \times E$ and
Stage~2 (level~$2$) restricting to the elliptic reference curve
$E_\tau$; the Borcherds product $\Phi_N$ supplies the BKM character
whose elliptic specialisation matches the $E_\tau$-collision residue
(level~$5$ scalar shadow under
$\kappa_{\mathrm{BKM}}(\Phi_N)=c_N(0)/2$).
\end{remark}

\section{Seven faces of $r(z)$ at the K3 crown}
\label{sec:g1sf-supplement}

\begin{remark}[Seven faces of $r(z)$ at the K3 crown]
\label{rem:g1sf-seven-faces}
On the Mukai-Heisenberg comparison lane, the genus-$1$ collision
kernel has a scalar K3 conductor
$K^{\kappa_{\mathrm{ch}}}(\mathcal{H}_{\mathrm{Muk}}(K3)) = 2\,c_+(\mathrm{Mukai}(K3)) = 8$.
Three scalar readouts of this $8$ are compatible with Bruinier
$2002$ Proposition~$5.1$ Heegner Chern-class reciprocity:
(a) Mukai doubling $2 c_+ = 8$;
(b) order of monodromy of $\mathcal{L}^{\Delta_5}$ around Humbert divisor $H_1 \subset \overline{\mathcal{A}_2}$;
(c) Lusztig specialisation $\ell = 8$ of the Hall--Drinfeld double at $\zeta^8 = 1$ with $\hbar^2 = -1/8$.
This is a scalar conductor statement: the Mukai Heisenberg algebra,
the Humbert divisor monodromy, and the Lusztig root-of-unity
Hall--Drinfeld double remain three full objects with a shared scalar
readout. Vol~I
Theorem~C uses the resulting scalar bucket
$\{0, 8, 13, 250/3, 98/3\}$.
\end{remark}

\begin{remark}[Siegel-elliptic comparison kernel]
\label{rem:g1sf-R-sieg}
At $K3\times E$, a genus-$2$ comparison target is the
Siegel-elliptic dynamical kernel
$R_{\mathrm{Sieg,dyn}}(u-v;\tau,\rho,z)$ built from the genus-$2$
theta function
$\theta[\alpha/\beta](w;\tau,z,\rho)$ on the abelian surface with
period matrix
$\bigl(\begin{smallmatrix}\tau & z\\ z & \rho\end{smallmatrix}\bigr)
\in\mathbb H_2$, using the Pasol--Zagier $2013$ extension of the
Kronecker--Eisenstein series. On the paramodular locus $K(1)$ this
kernel is the natural genus-$2$ dynamical-YBE comparison object. Its
restriction to the genus-$1$ collision residue is the comparison
obligation left outside Theorem~\ref{thm:g1sf-master}.
\end{remark}

\section{Seven K3-genus-$1$ comparison kernels}
\label{sec:g1sf-supplement-ii}

For an elliptic K3 with $24$ Kodaira $I_1$ fibres, the Vol~I
genus-$1$ kernel has K3 comparison maps on the smooth elliptic fibre,
near the nodal fibres, and on the Humbert boundary. The following
proposition records the polar kernels and scalar conductor classes
seen by those maps.

\begin{proposition}[K3-genus-$1$ comparison kernels]
\label{prop:g1sf-seven-K3-faces}
\index{seven faces!K3-genus-1|textbf}
\index{R-matrix@$R$-matrix!K3 catalogue|textbf}
\ClaimStatusConditional{} \textup{(}comparison maps; individual
external kernels:
Belavin $1981$; Harvey--Moore $1996$; Borcherds $1998$;
Gritsenko--Nikulin $1997$; Pasol--Zagier $2013$; Eguchi--Ooguri--Tachikawa
$2011$\textup{)}
\textbf{Type signature.} (Cross-volume Open$\leftrightarrow$CY at
level~$2$ for the bar/twisting kernel and level~$5$ for the K3
scalar conductor; cross-volume bridge to Vol~III at
$(\Sigma_{d-1}, C) = (\Sigma_2, E_\tau)$ via
$\Phi_3^{(\Sigma_2,E_\tau)} = \mathrm{Sp}^{\mathrm{ch}}_{\Sigma_2,E_\tau}\circ
\Phi_3^{\mathrm{FA}}$; hypothesis package: K3 elliptic-fibre datum
$(E_\tau,\text{Kodaira},\widetilde{M}_{24})$; comparison maps for
each kernel as cited; Mukai pairing on $H^{\mathrm{ev}}(K3,\Z)$.)
At genus $1$ with the K3-modular datum $(E_\tau, \text{Kodaira fibre
structure}, \widetilde{M}_{24})$, the following seven comparison
kernels have the stated polar or scalar readouts:
\begin{enumerate}
\item \textbf{(Rational face.)} On the smooth cover
$E_\tau^{\mathrm{sm}}\setminus\{24\}$, the Belavin--Drinfeld rational
$r$-matrix has polar part $r^{\mathrm{rat}}(z)=\Omega/z$, the
base-affine collision pole of the Heisenberg stalks.
\item \textbf{(Trigonometric face.)} Near each Kodaira node
$p_i \in E^{\mathrm{nod}}_{24}$, the nearby-cycles $\mathcal{D}$-module
has trigonometric Belavin--Drinfeld polar kernel $r^{\mathrm{trig}}(z)$
with spectral parameter $q_i = e^{2\pi i\tau_i}$ capturing the
$I_1$-monodromy.
\item \textbf{(Elliptic Belavin face.)} Away from nodes, the
Belavin elliptic kernel has Cartan polar part
$\xi_\tau(z)\,\Omega_{\mathfrak h}$ and root-channel theta ratios
\eqref{eq:g1sf-phi-pm}; this is the K3 restriction of Face~$5$.
\item \textbf{(Mukai-twisted face.)} The Mukai enhancement
$\mathcal{H}_{\mathrm{Muk}}(K3)=H^{\mathrm{ev}}(K3,\bZ)$ contributes a
Serre-functor cocycle term to the elliptic kernel,
\[
 r^{\mathrm{Muk}}(z; \tau) \;=\; r^{\mathrm{ell}}(z;\tau) + \mathrm{Serre}^{\widetilde{M}_{24}}(z),
\]
whenever the Mukai comparison functor to
$\mathrm{End}(\mathcal{H}_{\mathrm{Muk}})$ is fixed. Its scalar
conductor readout is $K^{\kappa_{\mathrm{ch}}}=8$.
\item \textbf{(Harvey--Moore face.)} The $1$-loop heterotic
threshold integrand on $K3\times T^2$
\cite{HarveyMoore1996} has logarithmic collision kernel
\[
 r^{\mathrm{HM}}(z;\tau)
 \;=\;
 -\,\frac{\theta_1'(z;\tau)}{\theta_1(z;\tau)}\,\Omega_{\Lambda^{3,2}},
\]
with $\Lambda^{3,2}$ the Narain lattice carrying the $24$ root-lattice
data. The comparison with Bruinier's Heegner Chern class is through
the scalar divisor class.
\item \textbf{(Borcherds-product face.)} The singular-theta lift of
the K3 elliptic genus $\phi_{0,1}$ by Borcherds 1998 has logarithmic
derivative
\[
 r^{\mathrm{Borch}}(z;\tau)
 \;=\;
 \partial_z \log\,\Delta_5(z,\tau,\rho)\bigl|_{\rho\to i\infty}
 \;=\;
 -\sum_{\alpha>0,\,\alpha^2\geq 0} m(\alpha)\,
 \frac{e^{2\pi i\alpha\cdot z}}{1 - e^{2\pi i\alpha\cdot z}}\,\Omega_\alpha,
\]
with multiplicities $m(\alpha)$ given by the Borcherds product.
\item \textbf{(Siegel-dynamical face.)} The
Siegel-elliptic dynamical kernel
$R_{\mathrm{Sieg,dyn}}(u-v; \tau,\rho,z)$ of
Remark~\ref{rem:g1sf-R-sieg} restricts, in the limit
$\rho\to i\infty$, to the genus-$1$ elliptic kernel together with the
Pasol--Zagier dynamical correction.
\end{enumerate}
At the Humbert divisor $H_1\cup H_4\subset\overline{\mathcal{A}_2}$,
the K3-specific comparison maps have the scalar $\ell=8$ conductor
readout of Remark~\ref{rem:g1sf-seven-faces}.
\end{proposition}

\begin{proof}
The proof is a comparison of kernels after applying the stated
extraction maps. The rational, trigonometric, and elliptic polar
kernels are the Belavin--Drinfeld rational/trigonometric/elliptic
classification \cite{Belavin81,BelavinDrinfeld82}, with the elliptic
Cartan term written in the prime-form convention of
\eqref{eq:g1sf-xi-kernel}. The Mukai term uses the CY-$2$ Serre shift
on $D^b(K3)$ and contributes only after choosing the Mukai comparison
functor to $\mathrm{End}(\mathcal{H}_{\mathrm{Muk}})$. The
Harvey--Moore kernel is the logarithmic derivative of the heterotic
threshold integrand \cite{HarveyMoore1996}; Bruinier's Heegner
Chern-class identity supplies the scalar divisor comparison. The
Borcherds-product kernel is the logarithmic derivative of the
singular-theta lift; Borcherds 1998 and Gritsenko--Nikulin 1997 give
the product and additive lift presentations of $\Delta_5$. The
Siegel-dynamical kernel is the Pasol--Zagier Kronecker--Eisenstein
extension on the genus-$2$ period matrix; restricting
$\rho\to i\infty$ gives the stated genus-$1$ kernel plus dynamical
correction. Bruinier $2002$ Proposition~$5.1$ gives the Humbert
scalar class; the Lusztig root-of-unity readout is the imposed
Hall--Drinfeld specialization $\zeta^8=1$.
\end{proof}

\begin{remark}[Faces vs.\ specialisations; naming discipline]
\label{rem:g1sf-faces-vs-specializations}
The seven K3-genus-$1$ entries above are comparison kernels, not seven
independent $r$-matrices. The common datum is the binary polar
collision kernel on the elliptic fibre. The Mukai, Harvey--Moore,
Borcherds-product, and Siegel-dynamical rows add scalar or dynamical
data visible only after the corresponding comparison map is chosen.
\end{remark}

\begin{remark}[K3 comparison and BV obstruction]
\label{rem:g1sf-cross-ref}
\ClaimStatusHeuristic{}
The expected relation between the vanishing of the $2$-loop BV obstruction class
(Proposition~\ref{prop:bvbrst-2loop-obstruction}) and the
seven K3-genus-$1$ face catalogue of
Proposition~\ref{prop:g1sf-seven-K3-faces} is the following numerical
matching problem: the theta-graph contribution carries the factor
$\chi(K3)=24$, while the nodal elliptic-fibre comparison sums the
$24$ Kodaira $I_1$ local trigonometric kernels. A proof would require
identifying this local kernel sum with the BV graph integral in the
same regularisation scheme.
\end{remark}
